\documentclass[uplatex, dvipdfmx, 11pt, a4paper, report]{jsbook}
\usepackage{float}
\usepackage{graphicx}
\usepackage{amsmath, amssymb}
\usepackage{url}
\usepackage{bm}
\usepackage{hyperref}
\usepackage{longtable}
\usepackage{subcaption}
\usepackage{listings}
\usepackage[dvipdfmx]{xcolor}

% コード表示のデザイン設定
\lstset{
  basicstyle=\ttfamily\footnotesize,
  commentstyle=\color{green!60!black},
  keywordstyle=\color{blue},
  stringstyle=\color{red},
  frame=single,
  numbers=left,
  numberstyle=\tiny\color{gray},
  breaklines=true,
  captionpos=b
}

\begin{document}


% 表紙

\begin{titlepage}
  \centering
  \vspace*{30mm} 
  
  {\huge 令和7年度 卒業論文 \par}
  
  \vspace{20mm}
  
  {\Huge ウェブ小説のメタデータに基づく\\連載放棄（エタり）の予測 \par}
  
  \vfill
  
  \begin{tabular}{rl}
      \Large 学籍番号 & \Large 2210306 \\[15pt]
      \Large 氏名     & \Large 佐藤 玲瑠 \\[15pt]
      \Large 学科     & \Large Ⅰ類メディア情報学プログラム \\[15pt]
      \Large 指導教員 & \Large 久野 雅樹 教授
  \end{tabular}
  
  \vspace{40mm} 
  
  {\Large 提出日 \today}
  
  \vspace{20mm}
\end{titlepage}

\newpage
\thispagestyle{empty}
\mbox{}
\newpage

% 目次
\setcounter{page}{1}
\tableofcontents


% 第1章 はじめに

\chapter{はじめに}

\section{研究の背景}

\subsection{ウェブ小説市場の拡大とUGCの隆盛}
近年，インターネット技術の発展とスマートフォンの普及に伴い，誰もが手軽に創作物を発表できるUGC（User Generated Content）プラットフォームが急速に成長している．
特に日本国内においては，ウェブ小説投稿サイト『小説家になろう』\footnote{\url{https://syosetu.com/}}を筆頭に，多数の投稿サイトが巨大な市場を形成している．
これらのサイトでは数多くの作品が日々投稿され，人気作品が書籍化，アニメ化，ゲーム化などのメディアミックス展開へと繋がるケースも一般化している．
ウェブ小説は今や，日本のエンターテインメント産業における重要な源泉となっており，その影響力は年々増大している．

\subsection{「エタり」の問題と読者の課題}
このように膨大な作品が日々投稿される一方で，完結することなく長期にわたり更新が停止してしまう作品も数多く存在する．
ウェブ小説コミュニティにおいては，未完結のまま更新が途絶える現象は俗に「エタる\footnote{Eternal（永遠）が語源とされている}」と呼称される．
『小説家になろう』のシステム上では，一定期間更新がない作品には「この作品は未完結のまま1年以上更新がありません」といった警告文が表示される（図\ref{fig:warning}参照）．

\begin{figure}[H]
  \centering
  \fbox{\includegraphics[width=12cm]{images/screenshot-novels-warning.png}}
  \caption{長期間更新が停止した作品に表示される警告文の例}
  \label{fig:warning}
\end{figure}

読者にとって，時間をかけて読み進めた物語が未完のまま放置されることは大きな課題である．
膨大な作品群の中から，自分の好みに合い，かつ完結まで読み通せる作品を見つけ出すことが困難な状況にある．

\subsection{執筆者側の課題と創作の持続可能性}
「エタり」は読者だけの問題ではない．執筆者にとっても，モチベーションの維持は困難な課題である．
十分なPVや感想といったフィードバックが得られないまま執筆を続けることは精神的な負担が大きく，結果として更新停止に至るケースが後を絶たない．
もし事前に「どのような要素が更新停止につながりやすいか」を客観的に把握できれば，執筆者は投稿スタイルや内容を改善し，連載放棄を回避するための指針を得ることができると考えられる．

\section{現状の課題と分析の動機}

\subsection{既存指標の限界}
現在，多くのウェブ小説サイトでは「ランキング」や「ブックマーク数」などが作品選びの指標となっている．
しかし，これらは「現時点での人気」を測るものであっても，「将来的に連載放棄されないかどうか」を保証するものではない．
読者が安心して作品を読み始めるためには，人気とは別の軸である連載放棄リスクを予測する新たな指標が必要とされている．

\subsection{予備調査による知見}
本研究に先立ち，著者が先行して収集したデータセットを用いて予備的な分析を行ったところ，興味深い事実が明らかになった．
例えば，あらすじに「完結」という単語が含まれている作品であっても，実際には未完結のまま放置されているケースが多数確認されたのである．
これは，「第一部完結」や「前作は完結済み」といった文脈で使用されているためであり，単純なキーワードマッチングでは更新停止リスクを正しく検知できないことを示唆している．
したがって，高性能な予測を行うためには，数値データの分析だけでなく文脈レベルでの解析が不可欠である．

\section{本研究の目的とアプローチ}

\subsection{研究の目的}
本研究の目的は，ウェブ小説のメタデータ（数値情報）およびあらすじ（テキスト情報）を用いて，その作品が将来的に「エタる」か否かを予測するモデルを構築することである．
また，構築したモデルを用いてリスク要因を可視化・分析することで，読者の作品選びを支援するとともに，執筆者の連載放棄回避を支援する知見を得ることを目指す．

\subsection{採用する手法}
本研究では，予測性能と解釈性のバランスを考慮し，多角的な視点から以下の手法を採用し比較検証を行う．

\begin{enumerate}
  \item \textbf{数値データに基づく分析}:
  投稿間隔，文字数，ブックマーク数などのメタデータに対し，LightGBMやロジスティック回帰などの機械学習モデルを適用する．
  \item \textbf{あらすじテキストに基づく分析}:
  対象作品のあらすじに対し，TF-IDF，Bi-LSTM，BERTなどの自然言語処理技術を適用し，文脈や意味的特徴に基づいた予測を試みる．
\end{enumerate}

\section{本論文の構成}
本論文は全6章から構成される．
第2章では，ウェブ小説の分析や，テキスト分類に関する関連研究について概観する．
第3章では，本研究で使用するデータセットの収集方法，前処理の詳細，および提案モデルの構築手法について述べる．
第4章では，構築した各モデルを用いた評価実験の結果を示し，性能比較を行う．
第5章では，得られた結果をもとに，あらすじ内のどのような表現が更新停止リスクに影響しているかを考察する．
最後に第6章で本研究の成果を総括し，今後の課題について述べる．



% 第2章 関連研究


\chapter{関連研究}

本章では，ウェブ小説を対象とした既存の分析研究について概観する．特に，ランキング予測，トレンド分析，多様性の変化，および言語的特徴の抽出に取り組んだ先行研究を取り上げ，それぞれの成果と課題について述べる．その上で，本研究の独自性と立ち位置を明確にする．

\section{清水らの研究（集合知に基づくランキング予測）}
ウェブ小説における最大の関心事は「どの作品が人気を獲得できるか」であるため，ランキング予測に関する研究は以前より精力的に行われている．
清水ら\cite{shimizu}は，作品に対する読者のブックマーク行動を「集合知」として捉え，ウェブページの評価アルゴリズムであるHITSを応用して将来のランキング上位作品を予測する手法を提案した．

彼らは，単純なブックマーク数だけでなく，「誰がブックマークしたか」という評価者の質を考慮するために，図\ref{fig:shimizu_graph}に示すような読者と小説の二部グラフ（ネットワーク構造）を構築した．
この構造を用いて，多くの良作をブックマークしている読者を「目利き（Hub）」として評価し，その目利きが選んだ作品のスコアを高く算出することで，従来のランキングよりも早期にヒット作を発見できることを示している．
彼らの研究は，メタデータやユーザー行動のネットワーク分析が作品の将来性を予測する上で強力な指標になることを示唆しているが，あくまで「商業的なヒット」を予測するものであり，作品が途中で放棄されるかどうかという「連載放棄」のリスクについては考慮されていない．

\clearpage
\begin{figure}[H]
  \centering
  \includegraphics[width=0.5\linewidth]{images/research-shimizu.png}
  \caption{ブックマークによる読者と小説の二部グラフ（清水ら\cite{shimizu}より引用）}
  \label{fig:shimizu_graph}
\end{figure}

\section{堺らの研究（流行語の変遷分析）}
堺ら\cite{sakai}は，「小説家になろう」のあらすじデータから流行語を抽出し，年次ごとのジャンルの変遷を可視化している．
彼らは，単なる単語の出現頻度ではなく，類似単語の出現頻度も考慮した独自の指標（$new\_tf$）を提案している．これは，例えば「オリンピック」と「五輪」のように，同一の概念を指す単語や表記揺れを，共起関係に基づく類似度を用いて統合的に評価する手法である．これにより，単一の単語頻度のみでは捉えきれない，実質的なトピックの流行度を算出している．

表\ref{tab:sakai_comparison}は，この手法に基づき算出された，2010年と2019年の10月における重要単語の上位を示したものである．
2010年時点では「青年」「少年」といった一般的な語彙が上位を占めているのに対し，2019年では「転移」が1位，「異（異世界）」が2位と，いわゆる「異世界転生」ジャンルを示唆する単語が圧倒的な頻度で出現している．特に「転移」は2010年時点では10位であったが，9年間で首位にまで上昇しており，このジャンルの爆発的な流行を裏付けている．
なお，表中に見られる「世界と日本」などの一部の不自然な単語については，形態素解析処理の誤りによるものであると報告されている．

\clearpage
\begin{table}[H]
  \centering
  \caption{類似単語の出現頻度比較（堺ら\cite{sakai}より抜粋）}
  \label{tab:sakai_comparison}
  \begin{tabular}{c|lr|lr} \hline
    Rank & \multicolumn{2}{c|}{2010年10月} & \multicolumn{2}{c}{2019年10月} \\
    & 単語 & 出現頻度 & 単語 & 出現頻度 \\ \hline
    1 & 青年 & 2058.2 & 転移 & 14877.2 \\
    2 & こと & 1911.7 & 異 & 14272.7 \\
    3 & 事 & 1861.9 & 別世界 & 11702.4 \\
    4 & 中 & 1823.5 & 世界と日本 & 11687.4 \\
    5 & 少年 & 1670.0 & 世界文化 & 11066.4 \\
    6 & 同じ星 & 1638.8 & 新しい世界 & 10389.3 \\
    7 & 彼女 & 1631.1 & 世界 & 10291.8 \\
    8 & 辰原 & 1572.7 & 不思議な世界 & 8551.6 \\
    9 & お付 & 1569.6 & 事 & 8497.1 \\
    10 & 転移 & 1566.5 & 現実世界 & 8449.5 \\ \hline
  \end{tabular}
\end{table}

\section{本田らの研究（キーワードの多様性分析）}
本田ら\cite{honda}は，作品に付与されたキーワードのエントロピーや類似度を分析し，ウェブ小説コンテンツの画一化について検証を行った．
図\ref{fig:honda_graph}は，彼らの研究によるキーワード間の類似度（平均SumCos値）の年次推移を示している．
グラフでは2009年頃に一時的な減少が見られるものの，全体としては緩やかな上昇傾向が続いている．
これは，多くの作品で共通したキーワードが使用される頻度が高まったことを意味しており，投稿される作品の多様性が減少し，特定のヒットパターンへの収斂が進んでいることを示唆している．

\begin{figure}[H]
  \centering
  \includegraphics[width=0.8\linewidth]{images/research-honda.png}
  \caption{平均SumCos値の年次推移（本田ら\cite{honda}より引用）}
  \label{fig:honda_graph}
\end{figure}

\section{黄の研究（ビッグデータとしての言語特徴解析）}
黄\cite{huang}は，ウェブ小説をビッグデータとして捉え，一般文芸と比較した際の語彙特性や，人気作品に共通する言語的特徴を詳細に解析している．
同氏は47万件に及ぶあらすじデータを対象に，分散表現を用いたトピックモデルであるTop2Vecを適用した．初期抽出された2,306のトピックに対し，階層的トピック削減（Hierarchical topic reduction）を行うことで，ウェブ小説を構成する主要な100トピックへの集約を行っている．

表\ref{tab:huang_topic_select}は，抽出されたトピックの一例である．
黄は，Top2Vecによって出力されたキーワード群と，実際にそのトピックに分類された作品を確認することで各トピックの意味を解釈している．
解析の結果，最も大きなサイズを持つトピック群の一つ（表中のTopic 2等）が，「王太子」「婚約」「破棄」といった単語で構成されていることが判明した．これは，「宮廷における婚約破棄」という特定のプロットが，ウェブ小説において最も支配的なトレンドであることを定量的に裏付けている．
その他にも，「スキル」「ステータス」（Topic 1）といったゲーム用語や，「書籍化」（Topic 12）といったメタな語彙が独立したトピックとして抽出されており，ウェブ小説が読者のニーズに合わせて高度にテンプレート化された言語空間であることが示唆されている．

\begin{table}[H]
  \centering
  \caption{ウェブ小説あらすじにおける特徴的なトピック抽出例（黄\cite{huang}より抜粋）}
  \label{tab:huang_topic_select}
  \begin{tabular}{cll} \hline
    Topic ID & 主な構成語（上位抜粋） & 解釈されるジャンル・特徴 \\ \hline
    Topic 1 & スキル, ステータス, レベル, アイテム, 装備 & ゲーム的要素（RPG） \\
    Topic 2 & 王子, 貴族, 相手, 婚約, 立場, 彼女, 令嬢 & 悪役令嬢・宮廷恋愛（婚約破棄） \\
    Topic 12 & 書籍, 発売, 買い, web, 完結, コミカライズ & 書籍化・商業展開 \\
    Topic 83 & 大地, もふもふ, 恵み, 尊い, 成長, 癒やし & 癒やし・スローライフ \\
    Topic 86 & vrmmo, エーミール, 赤井, ゲーム, リアル & VRMMO（仮想現実） \\ \hline
  \end{tabular}
\end{table}

\section{本研究の立ち位置}
以上の関連研究により，ウェブ小説における「トレンドの変遷」「多様性の減少」「特有の言語表現」については多くの知見が得られている．しかし，これらの研究はいずれもマクロな傾向分析や，人気作品の特徴抽出に留まっている．
個々の作品に着目し，「その作品が連載放棄（エタり）に至るか否か」というリスクを予測する研究は，本研究の調査の範囲では十分になされていない．



% 第3章 実験設定および分析手法


\chapter{実験設定および分析手法}

本章では，まず本研究の研究対象である「小説家になろう」の概要とジャンル構造について述べる．続いて，分析に使用するデータセットの構築手順，および適用する数値解析・自然言語処理の具体的な手法について説明する．
なお，実験はPython 3.11環境で行った．使用した主要なライブラリのバージョンおよび詳細な構築環境については，付録\ref{app:environment}に示す．

\section{小説投稿サイト「小説家になろう」}

\subsection{サイトの概要と規模}
「小説家になろう」は，株式会社ヒナプロジェクト\footnote{\url{https://hinaproject.co.jp/}}が運営する日本最大級の小説投稿サイトである．
ユーザーは無料で小説を投稿・閲覧することができ，作品の人気に応じて商業出版やメディアミックスが行われるエコシステムが確立されている．
2026年1月時点での公開データによると，掲載作品数は約140万作品，総エピソード数は約1,880万話，総文字数は約570億文字に達しており，極めて大規模なテキストデータが蓄積されているプラットフォームである．

\subsection{ジャンル分類の定義}
同サイトでは，作品の投稿時に作者自身がジャンルを選択するシステムが採用されており，読者はこのジャンル区分を利用して好みの作品を検索する．
公式サイトの定義\cite{narou_help}に基づき，各ジャンルの特徴を以下に示す．

\subsubsection*{大ジャンル：恋愛}
恋愛を主題とした作品群である．
\begin{description}
    \item[異世界] 異世界を主な舞台とし，恋愛を描く作品．
    \item[現実世界] 現実に近しい世界を舞台とし，恋愛を描く作品．
\end{description}

\subsubsection*{大ジャンル：ファンタジー}
架空の世界や魔法などの非現実的な要素を含む作品群である．
\begin{description}
    \item[ハイファンタジー] 現実世界とは異なる，完全に架空の世界を主な舞台とした作品．
    \item[ローファンタジー] 現実世界に近しい世界観の中に，魔法や異能などのファンタジー要素を取り入れた作品．
\end{description}

\subsubsection*{大ジャンル：文芸}
文学性や特定のテーマ性を重視した作品群である．
\begin{description}
    \item[純文学] 娯楽性よりも芸術性や文学的な表現に重きを置いた作品．
    \item[ヒューマンドラマ] 人と人との交流，人の一生，人間らしさ等を主題とした作品．
    \item[歴史] 過去の時代を舞台にした作品．時代小説のほか，タイムトラベルものもここに含まれる．
    \item[推理] 事件等の謎解きを行う過程を主体とした作品（ミステリー）．
    \item[ホラー] 読者に恐怖感を与えることを主題とした作品．
    \item[アクション] 戦闘描写やアクションシーンを主体とした作品．
    \item[コメディー] 読者を笑わせることを主題とした作品（ギャグ）．
\end{description}

\subsubsection*{大ジャンル：SF}
科学的な空想要素を含む作品群である．
\begin{description}
    \item[VRゲーム] VR（仮想現実）技術を利用したゲーム世界が舞台となる作品．
    \item[宇宙] 宇宙空間や他の惑星を舞台とした作品．
    \item[空想科学] 実在・非実在を問わず，科学技術や理論の要素を含む作品．
    \item[パニック] 天災，汚染，大事故，侵略等の危機的状況下を舞台とした作品．
\end{description}

\subsubsection*{大ジャンル：その他}
上記の大ジャンルに分類されない作品群である．
\begin{description}
    \item[童話] 幼年，児童に向けた内容の作品．
    \item[詩] 言葉の響きやリズムを用いて想いを表現する作品．
    \item[エッセイ] 個人的観点により思想や物事を論じている作品．
    \item[リプレイ] TRPG（テーブルトークRPG）のプレイ経過を文章形式で記録したもの．
    \item[その他] 上記のいずれのジャンルにも該当しない作品．
\end{description}

\subsection{「ノンジャンル」について}
本研究で収集したデータの中には，ジャンルが「ノンジャンル」と設定されている作品が存在する．これは現在の投稿システムでは選択できない区分であり，特殊な経緯を持つデータである．
調査の結果，「ノンジャンル」は2016年に行われたサイト内のジャンル再編に伴い発生したカテゴリであることが判明した．当時，ジャンル情報の更新が行われなかった作品に対し，運営側が一時的な措置として「ノンジャンル」を一括で割り当てた経緯がある．
現在もこのカテゴリに残っている作品は，2016年以降，作者による情報の更新や修正が行われていないことを意味する．
したがって，本研究において「ノンジャンル」に分類される作品は，主に2016年の再編以前に投稿された古い作品群であり，長期間にわたりメタデータの管理（ジャンル設定の更新など）が行われていないデータとして扱う．

\section{使用する機械学習モデルの概要}
本節では，本研究の分析において使用する主要な機械学習アルゴリズムおよび自然言語処理モデルの理論的背景について述べる．

\subsection{ロジスティック回帰 (Logistic Regression)}
ロジスティック回帰は，線形回帰モデルを分類問題に応用した一般化線形モデルの一種である．
入力特徴量の線形結合をシグモイド関数に通すことで，確率（0から1の値）を出力する．計算コストが低く解釈性が高いため，分類タスクにおけるベースラインとして広く用いられる．本研究では，他の複雑なモデルの性能を評価するための基準として使用する．

\subsection{サポートベクターマシン (SVM)}
サポートベクターマシン（Support Vector Machine: SVM）は，クラス間を分離する境界線（超平面）のうち，データ点との距離（マージン）が最大になるものを探索する手法である．
カーネル関数を用いることで，データを高次元空間へ写像し，線形分離不可能なデータに対しても非線形な境界線を引くことが可能である．

\subsection{決定木アンサンブル学習}
単一の決定木は解釈性が高いが過学習しやすいという欠点がある．これを克服するため，複数の決定木を組み合わせるアンサンブル学習が用いられる．

\subsubsection*{ランダムフォレスト (Random Forest)}
ランダムフォレストは，バギング（Bootstrap Aggregating）に基づく手法である．
訓練データからブートストラップサンプリング（復元抽出）を行って複数の決定木を独立に作成し，それらの多数決や平均によって最終的な予測を行う．

\subsubsection*{勾配ブースティング決定木 (GBDT)}
勾配ブースティング決定木（Gradient Boosting Decision Tree: GBDT）は，ブースティングに基づく手法である．
複数の決定木を逐次的に構築し，前の決定木の予測誤差（残差）を次の決定木が補正するように学習が進められる．
本研究では，GBDTの実装として以下の2つのライブラリを採用する．
\begin{itemize}
    \item \textbf{LightGBM}\cite{lightgbm}: Microsoftが開発した高速なGBDT実装である．Leaf-wise（葉単位）での木成長アルゴリズムを採用しており，大規模データに対しても高速に学習が可能である．
    \item \textbf{XGBoost}\cite{xgboost}: 拡張性と柔軟性に優れたGBDT実装である．正則化項を含んだ目的関数により過学習を抑制する機能が強力である．
\end{itemize}

\subsection{多層パーセプトロン (MLP)}
多層パーセプトロン（Multilayer Perceptron: MLP）は，入力層，隠れ層（中間層），出力層から構成される順伝播型ニューラルネットワークである．
隠れ層において活性化関数（ReLU等）を用いることで，入力データ間の複雑な非線形関係を学習することができる．

\subsection{自然言語処理モデル}
あらすじテキストの分析においては，以下の3つの手法を比較・検討する．

\subsubsection*{TF-IDF}
TF-IDFは，文書内の単語の重要度を数値化する統計的手法である．
単語の出現頻度（TF）と，その単語が登場する文書数の逆数（IDF）を掛け合わせることで，「特定の文書に頻出するが，全体では珍しい単語」に高い重みを与える．

\subsubsection*{双方向LSTM (Bi-LSTM)}
LSTM（Long Short-Term Memory）は，リカレントニューラルネットワーク（RNN）の一種であり，時系列データの長期依存関係を学習できる構造を持つ\cite{lstm}．
双方向LSTM（Bi-directional LSTM）は，過去から未来への順方向と，未来から過去への逆方向の2つのLSTMを組み合わせたモデルである．これにより，文脈を前後双方向から捉えることが可能となり，文章全体の意味理解において従来の単方向RNNよりも優れた性能を発揮する．

\subsubsection*{BERT}
BERT (Bidirectional Encoder Representations from Transformers) は，Googleが2018年に発表したTransformerベースの事前学習済み言語モデルである\cite{bert}．
大量のテキストデータを用いて「マスク化言語モデリング（Masked LM）」と「次文予測（NSP）」という2つのタスクで事前学習を行うことで，深い文脈理解能力を獲得している．
特定のタスクに合わせて少量のデータで追加学習（ファインチューニング）を行うことで，分類や質問応答など多様なNLPタスクにおいて高い性能を達成できることが知られている．

\section{データセットの構築}

\subsection{データ収集プロセス}
本研究では，ヒナプロジェクト社が提供する「なろう小説API\footnote{\url{https://dev.syosetu.com/}}」を利用し，網羅的なデータ収集を行った．
特定のランキングやジャンルに偏ることを防ぐため，2025年10月時点での登録ユーザーIDの2,814,253までを順次クエリとして送信し，各ユーザーが投稿している全小説のメタデータを取得し，生データセットを作成した．

\subsection{正解ラベルの定義とフィルタリング}
収集したデータに対し，予測の目的変数となる「エタり（長期更新停止）」のラベル\texttt{is\_eternal}を以下のように定義し，付与した．

\begin{itemize}
  \item \textbf{エタり (\texttt{is\_eternal} = 1)}: 
  作品設定が「未完結（連載中）」のまま，最終更新日から1年以上経過している作品．
  \item \textbf{非エタり (\texttt{is\_eternal} = 0)}: 
  作品設定が「完結済み」である作品．
\end{itemize}

また，本研究の目的は「連載の継続性」を予測することにあるため，分析に適さない以下の作品をデータセットから除外した．

\begin{enumerate}
  \item \textbf{連載中の作品}: 
  最終更新から1年未満かつ完結していない作品は，将来エタるか不明であるため除外した．
  \item \textbf{短編・読み切り作品}: 
  短編作品は投稿時点で完結扱いとなるため除外した．
\end{enumerate}

上記のフィルタリングの結果，最終的に分析に使用するデータセットは537,829件となった．
データセットに含まれる主なデータ項目を表\ref{tab:dataset_columns}に示す．

\begin{table}[H]
  \centering
  \caption{データセットに含まれる特徴量一覧}
  \label{tab:dataset_columns}
  \begin{tabular}{|l|l|} \hline
    分類 & 項目名 \\ \hline
    基本情報 & タイトル, Nコード, 作者ID, 作品ジャンル, 小説タイプ \\
    テキスト & あらすじ, キーワード \\
    日時情報 & 初回投稿日, 最終更新日 \\
    量的情報 & 投稿話数, 文字数, 読了時間, 会話率 \\
    評価指標 & 日間/週間/月間/四半期/年間ポイント, 総合評価ポイント \\
    反応数 & ブックマーク数, 感想数, レビュー数, 評価者数 \\
    目的変数 & is\_eternal (0 or 1) \\ \hline
  \end{tabular}
\end{table}

\section{数値データに基づく分析手法}

\subsection{特徴量の抽出とエンジニアリング}
収集したメタデータ（表\ref{tab:dataset_columns}参照）だけでなく，モデルが学習しやすいように既存の特徴量を用いて新たな特徴量を作成した．
本研究では，作品の時系列的な性質や執筆の継続性をより強く反映させるため，以下の3つの特徴量を新たに定義し，追加した．

\begin{enumerate}
  \item \textbf{更新頻度}: 
  1話投稿あたりの平均日数を更新頻度と定義した．
  \begin{equation}
    \text{更新頻度} = \frac{\text{最終更新日} - \text{初回投稿日}}{\text{投稿話数}} \quad (\text{単位: 日/話})
  \end{equation}

  \item \textbf{初回投稿年月指標 ($M_{start}$)}: 
  作品が投稿された時期を数値化するために，データセット内の最古の投稿日時（$D_{min}$）を基準とし，各作品の初回投稿日（$D_{start}$）までの経過月数を算出した．これにより，長期的なトレンド（昔の作品か，最近の作品か）をモデルに考慮させる．
  \begin{equation}
    M_{start} = 12 \times (Y_{start} - Y_{min}) + (m_{start} - m_{min})
  \end{equation}
  ここで，$Y$は年，$m$は月を表す．

  \item \textbf{連載経過日数 ($D_{duration}$)}: 
  その作品がどれだけの期間執筆されていたか（連載期間）を表す指標である．初回投稿日（$D_{start}$）から最終更新日（$D_{last}$）までの日数を算出した．
  \begin{equation}
    D_{duration} = D_{last} - D_{start} \quad (\text{単位: 日})
  \end{equation}
\end{enumerate}

最終的にモデルに入力する特徴量は，文字数，話数，ブックマーク数などの基本メタデータに加え，これらの派生特徴量を含めたものである．
分析は，これらの特徴量を用いて「全データセットを一括」および「ジャンル別」の2パターンで実施した．

\subsection{使用モデルと学習設定}
数値分類タスクにおけるベースラインおよび比較対象として，以下の7つの機械学習モデルを採用した．
\begin{enumerate}
  \item 勾配ブースティング決定木: LightGBM, XGBoost
  \item 決定木アンサンブル: ランダムフォレスト (Random Forest)
  \item 線形モデル: ロジスティック回帰 (Logistic Regression)
  \item その他: サポートベクターマシン (SVM), 多層パーセプトロン (MLP), k近傍法 (k-NN)
\end{enumerate}

「エタり」と「非エタり」のデータ数は大きく偏っているため，学習時には作品数が少ないクラスに合わせるアンダーサンプリングにより調整を行った．
モデルの評価には5分割交差検証を用い，汎化性能を測定した．
また，決定木ベースのモデル（LightGBM, Random Forest, XGBoost）については，特徴量重要度を算出し，どの変数が予測に寄与しているかを分析した．

\subsection{誤分類分析}
モデルの予測信頼度（確率）と実際の結果との乖離を分析するため，予測確率が80\%以上であるにも関わらず予測を外した事例を抽出した．
これらの事例について，特徴量の分布を可視化し，モデルが苦手とするパターンの特定を試みた．

\section{あらすじテキストに基づく分析手法}

あらすじのテキスト情報を用いた分析では，手法の複雑さに応じて以下の3段階のアプローチを行った．
いずれの手法においても，前処理としてテキスト内のURLを特殊トークン \texttt{<URL>} に置換し，数値分析と同様にアンダーサンプリングによるデータ均衡化を行った．

\subsection{TF-IDFと機械学習モデルによる分析}
ベースラインとして，BoW (Bag-of-Words) アプローチを用いた．
形態素解析エンジンMeCab\cite{mecab}を用いてテキストを分かち書きし，品詞を「名詞，動詞，形容詞，接頭詞」に限定して抽出した（非自立語，代名詞，数は除外）．
抽出した単語をTF-IDFベクトルに変換し，ロジスティック回帰等の分類器に入力した．
ハイパーパラメータの探索にはグリッドサーチを用いた．

\subsection{SentencePieceとBi-LSTMによる分析}
文脈を考慮した系列データとしての分析を行うため，サブワード分割を用いたニューラルネットワークモデルを構築した．
トークナイザにはSentencePiece\cite{sentencepiece}を採用し，語彙サイズは訓練データサイズの半分（最小5,000〜最大30,000）に設定した．
モデル構造には双方向LSTM (Bi-LSTM) を採用し，最大エポック数を10とした．学習時にはEarly Stoppingを適用し，過学習を防いだ．

\subsection{BERTによる分析}
文脈の深い意味理解を行うため，事前学習済みの大規模言語モデルBERTを用いた．
モデルには東北大学が公開している日本語モデル \texttt{tohoku-nlp/bert-large-japanese-v2} を採用した．
計算コストの観点から，学習データは各クラス最大5,000件，合計10,000件にサンプリングして使用した．

\subsection{LIMEによる判断根拠の可視化}
構築したモデルがテキスト中の「どのような単語」に着目してエタり判定を行っているかを明らかにするため，LIME (Local Interpretable Model-agnostic Explanations) を用いた分析を行った．
LIMEは，ブラックボックスな機械学習モデルに対し，局所的な近似を行うことで予測の根拠を提示する手法である．

本研究では，以下の手順でジャンルごとの重要単語を抽出した．
\begin{enumerate}
  \item \textbf{サンプリング}: 
  計算コストを考慮し，対象ジャンルのテストデータから「エタり」「完結」をそれぞれ同数ずつランダムにサンプリングした（全作品の場合は各5,000件，合計10,000件）．
  \item \textbf{局所的な説明の生成}: 
  サンプリングした各作品に対し，学習済みモデル（TF-IDF，Bi-LSTM，BERT）を用いてLIMEによる分析を行い，その作品の判定に寄与した単語とその寄与度（スコア）を算出した．
  \item \textbf{スコアの集計}: 
  全サンプリングデータについて単語ごとのスコアを合計し，最終的に「エタり判定に最も寄与した単語 Top15」および「完結判定に最も寄与した単語 Top15」を抽出した．
\end{enumerate}

なお，計算リソースの制約上，TF-IDFモデルについては全ジャンルで分析を実施したが，Bi-LSTMおよびBERTモデルについては「全作品（All）」データのみを対象として分析を行った．



% 第4章 実験結果


\chapter{実験結果}

本章では，第3章で構築したデータセットおよびモデルを用いた実験の結果について述べる．
まず，データセット全体の基礎的な統計情報を確認した後，数値データを用いた予測結果，あらすじテキストを用いた予測結果の順に詳細を示す．

\section{データセットの特性と基礎分析}

\subsection{ジャンルごとの連載放棄率（エタり率）}
作成したデータセットにおける「エタり率（連載放棄作品の割合）」をジャンルごとに算出した結果を表\ref{tab:eternal_rate}に示す．
なお，作品数が極端に少ない「未設定（4件）」は集計から除外した．

表を見ると，ジャンルによってエタり率に大きな開きがあることがわかる．
「童話」や「推理」，「ホラー」ジャンルでは過半数が連載放棄状態にあり，完結させることの難易度が極めて高いことが示唆される．
一方で，「アクション」や「VRゲーム」，「ノンジャンル」は相対的にエタり率が20\%未満と低く，これらのジャンルでは連載が継続されやすい，あるいは完結まで持っていきやすい傾向があると考えられる．

\clearpage
\begin{table}[H]
  \centering
  \caption{ジャンルごとのエタり率（全作品対象・降順）}
  \label{tab:eternal_rate}
  \begin{tabular}{lrr} \hline
    ジャンル & 作品総数 & エタり率 (\%) \\ \hline
    童話 & 3,098 & 61.91 \\
    推理 & 7,464 & 52.61 \\
    ホラー & 8,954 & 51.59 \\
    純文学 & 7,050 & 41.59 \\
    異世界（恋愛） & 48,842 & 40.55 \\
    歴史 & 6,442 & 40.07 \\
    詩 & 2,321 & 37.35 \\
    宇宙 & 3,117 & 36.67 \\
    現実世界（恋愛） & 46,262 & 35.14 \\
    ヒューマンドラマ & 36,913 & 34.82 \\
    空想科学 & 10,943 & 34.41 \\
    リプレイ & 492 & 32.72 \\
    エッセイ & 8,693 & 30.37 \\
    コメディー & 15,059 & 30.17 \\
    ローファンタジー & 47,058 & 26.94 \\
    パニック & 2,970 & 26.46 \\
    その他 & 15,898 & 25.90 \\
    ハイファンタジー & 122,213 & 21.28 \\
    アクション & 16,461 & 19.77 \\
    VRゲーム & 6,277 & 19.37 \\
    ノンジャンル & 121,298 & 16.63 \\ \hline
  \end{tabular}
\end{table}

\section{数値データに基づく分析結果}

投稿間隔，文字数，ブックマーク数などの数値データ（メタデータ）を入力とした場合の各モデルの予測性能を示す．
評価指標にはAccuracy（正解率），Precision（適合率），Recall（再現率），F1-Score（F1値）を用いた．

なお，本実験で得られた各ジャンルにおける「LightGBｍ，XGBoost，RandomForestそれぞれの特徴量重要度」および「誤分類グラフ」については，以下のオンライン補足資料\footnote{\url{https://sites.google.com/gl.cc.uec.ac.jp/sato-2025-research?usp=sharing}}に掲載している．

\subsection{全体データにおける予測性能}
データセット全体を対象にアンダーサンプリングを行い，両クラス合計221,424件に調整したデータセットに対する実験結果を表\ref{tab:result_all}に示す．
比較の結果，決定木ベースの勾配ブースティング手法である\textbf{XGBoost}および\textbf{LightGBM}が最も高い性能を示し，F1値で約0.78を記録した．
次いでランダムフォレスト，多層パーセプトロン（MLP）が続き，線形モデルであるロジスティック回帰は最も低い性能となった．
これは，ウェブ小説の「エタり」に関わる特徴量が線形分離不可能であり，複雑な相互作用を持っていることを示唆している．

\begin{table}[H]
  \centering
  \caption{数値データによる予測性能（全ジャンル統合）}
  \label{tab:result_all}
  \begin{tabular}{lrrrr} \hline
    Model & Accuracy & Precision & Recall & F1 Score \\ \hline
    XGBoost & \textbf{0.783} & 0.789 & 0.775 & \textbf{0.782} \\
    LightGBM & 0.782 & 0.788 & 0.773 & 0.780 \\
    RandomForest & 0.764 & 0.770 & 0.755 & 0.762 \\
    MLP (NeuralNet) & 0.754 & 0.749 & 0.765 & 0.756 \\
    SVM & 0.700 & 0.673 & 0.778 & 0.722 \\
    LogisticRegression & 0.658 & 0.625 & 0.787 & 0.697 \\ \hline
  \end{tabular}
\end{table}

\subsection{主要ジャンルごとの予測性能}
ウェブ小説の傾向はジャンルによって大きく異なるため，作品数の多い主要なジャンルについて個別にモデルを構築し，評価を行った．
ここでは代表的な例として「異世界（恋愛）」「現実世界（恋愛）」「ハイファンタジー」の結果を示す．
なお，その他のジャンルの詳細な結果については付録\ref{app:numeric_detail}に掲載する．

\subsubsection{異世界（恋愛）}
表\ref{tab:result_isekai_love}に「異世界（恋愛）」ジャンルの結果を示す．
このジャンルではXGBoostおよびLightGBMがF1値で0.87を超える極めて高い性能を記録した．
これは全体平均（0.78）と比較しても有意に高く，このジャンルの作品には連載放棄に至る特有の，あるいはパターン化された傾向が強く存在することがうかがえる．

\begin{table}[H]
  \centering
  \caption{予測性能：異世界（恋愛）}
  \label{tab:result_isekai_love}
  \begin{tabular}{lrrrr} \hline
    Model & Accuracy & Precision & Recall & F1 Score \\ \hline
    XGBoost & \textbf{0.875} & 0.880 & 0.868 & \textbf{0.874} \\
    LightGBM & 0.873 & 0.880 & 0.865 & 0.872 \\
    RandomForest & 0.846 & 0.853 & 0.837 & 0.845 \\
    MLP (NeuralNet) & 0.841 & 0.859 & 0.817 & 0.837 \\ \hline
  \end{tabular}
\end{table}

\subsubsection{現実世界（恋愛）}
表\ref{tab:result_reality_love}に「現実世界（恋愛）」の結果を示す．
同じ恋愛ジャンルであっても，異世界設定と比較すると性能は10ポイント近く低下しており，全体平均と同程度の0.78前後となった．
現実世界を舞台とした作品は設定の自由度が高く，画一的なパターンに当てはまりにくいことが要因の一つと考えられる．

\begin{table}[H]
  \centering
  \caption{予測性能：現実世界（恋愛）}
  \label{tab:result_reality_love}
  \begin{tabular}{lrrrr} \hline
    Model & Accuracy & Precision & Recall & F1 Score \\ \hline
    LightGBM & \textbf{0.784} & 0.790 & 0.773 & \textbf{0.781} \\
    XGBoost & 0.783 & 0.788 & 0.773 & 0.780 \\
    RandomForest & 0.766 & 0.773 & 0.753 & 0.763 \\ \hline
  \end{tabular}
\end{table}

\subsubsection{ハイファンタジー}
表\ref{tab:result_high_fantasy}に，最も作品数が多い「ハイファンタジー」の結果を示す．
このジャンルにおいてもLightGBMが最も良い性能を示し，F1値は約0.79であった．
データ数が最も豊富であるにも関わらず，異世界恋愛ほど性能が向上しなかった点は，ハイファンタジーというジャンルが内包する作品の多様性（王道冒険譚から復讐劇，スローライフまで多岐にわたる）を示唆している．

\begin{table}[H]
  \centering
  \caption{予測性能：ハイファンタジー}
  \label{tab:result_high_fantasy}
  \begin{tabular}{lrrrr} \hline
    Model & Accuracy & Precision & Recall & F1 Score \\ \hline
    LightGBM & \textbf{0.789} & 0.789 & 0.788 & \textbf{0.788} \\
    XGBoost & 0.788 & 0.789 & 0.787 & 0.788 \\
    RandomForest & 0.768 & 0.768 & 0.768 & 0.768 \\ \hline
  \end{tabular}
\end{table}

\subsection{特徴量の重要度分析}
構築した決定木ベースのモデル（XGBoost, LightGBM）において，予測に大きく寄与した特徴量を確認した．
全ジャンルおよび主要ジャンルを通して，共通して以下の特徴量が上位にランクインする傾向が見られた．

\begin{itemize}
  \item \textbf{連載経過日数}: 
  今回新たに追加した特徴量である．執筆期間が長い作品ほど，完結するかエタるかの傾向が明確に分かれるため，モデルが最も重視する指標の一つとなったと考えられる．
  \item \textbf{初回投稿年月}: 
  「なろう」のブームや時期による傾向（昔の作品は完結しやすい，最近はエタりやすい等）を捉えていると推測される．
  \item \textbf{投稿間隔（更新頻度）}: 
  更新が滞りがちな作品はエタりやすいという直感的な傾向を，モデルも定量的に捉えている．
  \item \textbf{ブックマーク数・評価ポイント}: 
  読者の反響が大きい作品は執筆モチベーションが維持されやすい等の相関があると考えられる．
\end{itemize}

特に「連載経過日数」と「初回投稿年月」の寄与度が非常に高いことは，ウェブ小説の連載放棄リスクを予測する上で，作品の内容（テキスト）だけでなく，「いつ，どのくらいの期間書いているか」という時間的なコンテキストが極めて重要であることを示唆している．

% 特徴量重要度
\begin{figure}[H]
  \centering
  % 上段：左
  \begin{subfigure}[b]{0.48\textwidth}
    \centering
    \includegraphics[width=\linewidth]{images/train_all-XGBoost-importance.png}
    \caption{XGBoost}
  \end{subfigure}
  \hfill
  % 上段：右
  \begin{subfigure}[b]{0.48\textwidth}
    \centering
    \includegraphics[width=\linewidth]{images/train_all-LightGBM-importance.png}
    \caption{LightGBM}
  \end{subfigure}
  
  \vspace{0.5cm}
  
  % 下段：中央
  \begin{subfigure}[b]{0.6\textwidth}
    \centering
    \includegraphics[width=\linewidth]{images/train_all-RandomForest-importance.png}
    \caption{RandomForest}
  \end{subfigure}
  
  \caption{全データセットにおけるモデル別特徴量重要度}
  \label{fig:feature_importance_all}
\end{figure}
  
  \vspace{0.5cm}
  
  % 異世界（恋愛）
\begin{figure}
  \centering
  \begin{subfigure}[b]{0.6\textwidth}
    \centering
    \includegraphics[width=\linewidth]{images/train_101-LightGBM-importance.png}
    \caption{異世界（恋愛）: LightGBM重要度}
    \label{fig:isekai_imp}
  \end{subfigure}
  
  \caption{全データセットおよび主要ジャンルにおける特徴量重要度}
  \label{fig:feature_importance}
\end{figure}

\subsection{誤分類（エラー）の傾向分析}
モデルが高い確信度を持って予測したにも関わらず外れた誤分類データについて，その密度分布を分析した．
その結果，ジャンルごとに特徴的なエラーの傾向が確認された．

% 全作品の誤分類傾向
\begin{figure}[H]
  \centering
  \begin{subfigure}[b]{0.48\textwidth}
    \centering
    \includegraphics[width=\linewidth]{images/train_all-classmiss1.png} 
    \caption{完結予測の誤り（偽陽性）}
    \label{fig:all_error_fp}
  \end{subfigure}
  \hfill
  \begin{subfigure}[b]{0.48\textwidth}
    \centering
    \includegraphics[width=\linewidth]{images/train_all-classmiss2.png} 
    \caption{エタり予測の誤り（偽陰性）}
    \label{fig:all_error_fn}
  \end{subfigure}
  \caption{全データセットにおける誤分類事例の特徴量分布}
  \label{fig:all_error_dist}
\end{figure}

\subsubsection{投稿話数における特異なピーク}
全ジャンルおよび「純文学」「推理」ジャンルにおいて，投稿話数の本来の分布密度は平坦でなだらかであるにも関わらず，誤分類データの分布には特定の投稿話数帯に鋭いピークが確認された．

\begin{itemize}
    \item \textbf{全ジャンル（完結予測の誤り）}: 
    特定の話数で終わっている作品に対し，モデルが過剰に「エタり（あるいは完結）」と判定を誤る傾向がある．これは，短編から長編への過渡期にある中編作品（例えば20〜50話程度など）において，完結か連載放棄かの判別が難しい境界線が存在することを示唆している．
    \item \textbf{純文学・推理}: 
    これらのジャンルは，ストーリー構成が多様であり，一般的な「なろう系」の法則（長編化しやすい等）が当てはまらないケースが多い．そのため，特定の長さの作品に対してモデルがバイアスを持った予測をしてしまい，誤分類が集中したと考えられる．
\end{itemize}

図\ref{fig:compare_error}に，予測性能が高かった「異世界（恋愛）」と，誤分類に特徴が見られた「純文学」の比較を示す．
異世界（恋愛）では誤分類の分布がなだらかであるのに対し，純文学では特定の話数にピークが立っていることが視覚的にも確認できる．

% 純文学と異世界恋愛の比較
\begin{figure}[H]
  \centering
  % 左：純文学
  \begin{subfigure}[b]{0.48\textwidth}
    \centering
    \includegraphics[width=\linewidth]{images/train_301-classmiss1.png}
    \caption{純文学：投稿話数の誤分類}
    \label{fig:purelit_error}
  \end{subfigure}
  \hfill
  % 右：異世界恋愛
  \begin{subfigure}[b]{0.48\textwidth}
    \centering
    \includegraphics[width=\linewidth]{images/train_101-classmiss1.png}
    \caption{異世界（恋愛）：投稿話数の誤分類}
    \label{fig:isekai_error}
  \end{subfigure}
  \caption{ジャンルによる誤分類傾向の違い（投稿話数）}
  \label{fig:compare_error}
\end{figure}

\subsubsection{更新頻度とジャンルの特異性}
「リプレイ」ジャンル等の特定カテゴリでは，更新頻度（投稿間隔）の誤分類分布が，本来の分布とは異なる形状を示した．
これは，TRPGのリプレイなど特殊な形式の作品では，「不定期更新でも完結する（あるいはエタる）」というパターンが他ジャンルと異なり，一般的な更新間隔のルールが通用しにくいことが要因と考えられる．

\subsubsection{特徴量による分類不能なケース}
「その他」ジャンルにおける完結予測の誤分類分布において，文字数や読了時間などの特徴量が横一直線の分布（均一な密度）を示した．
特定のピークが存在しないということは，モデルが特定の特徴量に基づいて判断できていない，すなわち「ランダムに近い状態で間違えている」ことを意味する．
「その他」ジャンルには定型化できない多様な作品が混在しているため，数値データ（メタデータ）だけでは予測の手がかりが掴めず，分類が困難であったと推察される．

これらの分析から，数値データを用いた予測モデルは全体として高い性能を誇るものの，ジャンル特有の執筆スタイルや，中編規模の作品においては判断に迷いが生じやすいことが明らかとなった．

\section{あらすじテキストに基づく分析結果}

\subsection{TF-IDFを用いたテキスト分類の結果}
あらすじや本文などのテキストデータからTF-IDFベクトルを作成し，分類モデル（ロジスティック回帰）に適用した結果を述べる．
グリッドサーチにより探索したハイパーパラメータ（特徴量次元数 $N$，正則化パラメータ $C$）のうち，最も高い正解率を記録した組み合わせと，その時の性能を表\ref{tab:tfidf_summary}に示す．
また，全ジャンルにおける探索結果の詳細は付録\ref{app:tfidf_grid}に掲載する．

\begin{table}[H]
  \centering
  \caption{TF-IDFモデルにおけるジャンル別最高性能と最適パラメータ}
  \label{tab:tfidf_summary}
  \begin{tabular}{lrrr} \hline
    ジャンル & 特徴量次元数 ($N$) & 正則化 ($C$) & 最高性能 \\ \hline
    \textbf{全データ (All)} & 20,000 & 0.20 & 0.6919 \\ \hline
    リプレイ & 700 & 0.05 & 0.8154 \\
    異世界（恋愛） & 8,000 & 0.80 & 0.7393 \\
    ヒューマンドラマ & 18,000 & 0.50 & 0.6849 \\
    現実世界（恋愛） & 22,000 & 0.80 & 0.6843 \\
    推理 & 18,000 & 1.00 & 0.6848 \\
    童話 & 3,000 & 0.08 & 0.6780 \\
    コメディー & 18,000 & 1.00 & 0.6700 \\
    アクション & 1,000 & 1.00 & 0.6797 \\
    ホラー & 10,000 & 1.00 & 0.6736 \\
    ハイファンタジー & 20,000 & 0.50 & 0.6697 \\
    宇宙 & 1,000 & 1.20 & 0.6681 \\
    パニック & 1,000 & 0.05 & 0.6730 \\
    歴史 & 7,000 & 0.50 & 0.6554 \\
    空想科学 & 5,000 & 1.00 & 0.6516 \\
    詩 & 1,000 & 0.05 & 0.6513 \\
    エッセイ & 2,000 & 1.50 & 0.6468 \\
    ノンジャンル & 5,000 & 1.20 & 0.6430 \\
    ローファンタジー & 3,000 & 0.50 & 0.6417 \\
    VRゲーム & 800 & 0.08 & 0.6386 \\
    純文学 & 1,000 & 0.80 & 0.6266 \\
    その他 & 10,000 & 1.00 & 0.6286 \\ \hline
  \end{tabular}
\end{table}

表\ref{tab:tfidf_summary}の結果より，以下の傾向が確認できる．

\begin{itemize}
    \item \textbf{特定ジャンルの高さ}: 「リプレイ」ジャンルにおいて約81.5\%という極めて高い性能が得られた．これは，TRPGのリプレイ形式という特殊な文体や，システム特有の専門用語（「ダイス」「GM」など）が強く機能したためと考えられる．
    \item \textbf{恋愛ジャンルの優位性}: 「異世界（恋愛）」においても約74\%と比較的高い性能を記録した．これは先行研究にあるように，流行語や特定のキーワード（「悪役令嬢」「婚約破棄」等）が明確なトレンドとして存在するため，単語ベースのモデルでも捉えやすかったと推測される．
    \item \textbf{メタデータモデルとの比較}: 全体（All）の性能は約69\%に留まり，前節のメタデータを用いたモデル（XGBoostやLightGBMで達成した約78\%）と比較すると性能は劣る結果となった．これは，文体や単語の出現頻度だけでは，「エタるか否か」という連載放棄リスクを完全に予測するのは困難であることを示唆している．
\end{itemize}

\subsection{サブワード分割とBi-LSTMによる分析結果}
次に，文脈を考慮した系列モデルであるBi-LSTM（Bidirectional LSTM）を用いた分析結果を述べる．
テキストのトークン化にはSentencePiece（サブワード分割）を用い，語彙サイズは学習データの規模に応じて調整した．
各ジャンルにおける予測正解率およびF1値を表\ref{tab:bilstm_result}に示す．
ジャンル別の詳細な評価指標については付録\ref{app:bilstm_detail}に掲載する．

\begin{table}[H]
  \centering
  \caption{Bi-LSTMモデルのジャンル別予測性能}
  \label{tab:bilstm_result}
  \begin{tabular}{lrr} \hline
    ジャンル & Accuracy & Macro F1 \\ \hline
    \textbf{全データ (All)} & 0.71 & 0.71 \\ \hline
    異世界（恋愛） & 0.72 & 0.72 \\
    ヒューマンドラマ & 0.68 & 0.68 \\
    現実世界（恋愛） & 0.67 & 0.67 \\
    童話 & 0.67 & 0.67 \\
    アクション & 0.66 & 0.66 \\
    推理 & 0.66 & 0.66 \\
    ハイファンタジー & 0.66 & 0.66 \\
    ホラー & 0.64 & 0.64 \\
    ノンジャンル & 0.64 & 0.63 \\
    コメディー & 0.63 & 0.63 \\
    ローファンタジー & 0.63 & 0.63 \\
    歴史 & 0.62 & 0.62 \\
    パニック & 0.61 & 0.61 \\
    宇宙 & 0.60 & 0.60 \\
    空想科学 & 0.60 & 0.59 \\
    エッセイ & 0.60 & 0.59 \\
    その他 & 0.60 & 0.59 \\
    詩 & 0.60 & 0.57 \\
    VRゲーム & 0.59 & 0.59 \\
    純文学 & 0.58 & 0.58 \\
    リプレイ & 0.51 & 0.34 \\ \hline
  \end{tabular}
\end{table}

\clearpage

表\ref{tab:bilstm_result}の結果から，TF-IDFと比較して以下の特徴が見られた．

\begin{itemize}
    \item \textbf{全体性能の向上}: 
    全データ（All）におけるAccuracyは0.71となり，TF-IDFの0.69から約2ポイント向上した．これは，単語の出現有無だけでなく，単語の並び順（文脈）を考慮することで，より精緻な予測が可能になったためと考えられる．
    \item \textbf{小規模データにおける学習不全}: 
    特筆すべき点として，「リプレイ」ジャンルの正解率が0.51（F1値 0.34）と著しく低下した．内訳を確認すると，再現率が0.00となっており，連載放棄作品を1件も正しく検出できていなかった．
    TF-IDFでは0.8154の正解率が出ていたことから，データ数が極端に少ない（学習・テスト合わせて数百件程度）ジャンルにおいては，パラメータ数の多いニューラルネットワークモデルは学習が収束せず，単純なキーワードマッチングであるTF-IDFの方が有効に機能することが示唆された．
    \item \textbf{主要ジャンルの安定性}: 
    「異世界（恋愛）」や「ヒューマンドラマ」などの作品数が多い主要ジャンルでは，軒並み高い性能を維持しており，データ量が十分にあれば深層学習モデルが有効であることが確認できた．
\end{itemize}

また，各ジャンルにおける詳細な「サブワード分割モデルの学習曲線（Loss推移）」については，前述のオンライン補足資料ウェブサイトにて公開している．

\section{BERTモデルによる分析結果}

あらすじを入力データとし，BERTを用いたファインチューニングを行ったモデルの評価結果を述べる．

\subsection{全体および主要ジャンルの傾向}
全ジャンルを統合した評価結果は，正解率が0.717，F1値が0.742となった．
主要な人気ジャンルについて確認すると，データ数が最も多い「異世界（恋愛）」においては正解率0.743，F1値0.769と，全体平均を上回る安定した性能を示した．これは，当該ジャンルのあらすじに特有の定型的な表現やキーワード（「婚約破棄」「悪役令嬢」等）が含まれやすく，BERTがその特徴を有効に学習できたためと考えられる．
一方，「現実世界（恋愛）」の正解率は0.687となり，異世界（恋愛）と比較するとやや低い値にとどまった．これは，現実世界を舞台とする作品の方が設定や展開の自由度が高く，語彙の多様性が高いため予測難易度が上がったことが示唆される．

\subsection{ジャンルごとの詳細評価}
表\ref{tab:bert_all_score}に，全ジャンルにおける詳細な評価指標を示す（F1値降順）．

その他のジャンルにおける特筆すべき結果として，以下の点が挙げられる．
\begin{itemize}
    \item \textbf{リプレイの大幅な改善}:
    Bi-LSTMモデルでは学習が困難であった「リプレイ」ジャンルにおいて，BERTモデルは正解率0.723，再現率0.906という極めて高い数値を記録した．これは，事前学習モデルであるBERTが，会話文やログ形式に近いリプレイ独特の文脈を，少数のデータからでも捉えることに成功した結果である．
    
    \item \textbf{アクション・童話の高精度}:
    「アクション」（正解率0.734）や「童話」（正解率0.728）においても，全体平均以上の性能が得られた．特にアクションはF1値も0.743と高く，適合率と再現率のバランスが良い結果となった．
    
    \item \textbf{エッセイ・VRゲームの低迷}:
    性能が低かったジャンルとして「エッセイ」（正解率0.593）と「VRゲーム」（正解率0.596）が挙げられる．特にVRゲームは再現率こそ高い（0.868）ものの，適合率が0.561と低く，過剰に「エタる」と判定してしまう傾向が見られた．
\end{itemize}

% 全ジャンルの詳細表
\setlength\LTleft{0pt plus 1fill}
\setlength\LTright{0pt plus 1fill}

\begin{longtable}{lcccc}
\caption{BERTモデル ジャンル別評価結果（全体およびF1値降順）} \label{tab:bert_all_score} \\
\hline
Genre & Accuracy & F1-score & Precision & Recall \\ \hline
\endfirsthead

\multicolumn{5}{c}{{続・表 \ref{tab:bert_all_score}}} \\
\hline
Genre & Accuracy & F1-score & Precision & Recall \\ \hline
\endhead

\hline \multicolumn{5}{r}{\textit{Continued on next page}} \\
\endfoot

\hline
\endlastfoot

% 全体結果
\textbf{全体 (All)} & 0.717 & 0.742 & 0.682 & 0.812 \\ \hline
% 各ジャンル
異世界（恋愛） & 0.743 & 0.769 & 0.699 & 0.855 \\
リプレイ & 0.723 & 0.763 & 0.659 & 0.906 \\
童話 & 0.728 & 0.759 & 0.680 & 0.860 \\
アクション & 0.734 & 0.743 & 0.719 & 0.770 \\
ホラー & 0.724 & 0.736 & 0.706 & 0.769 \\
コメディー & 0.715 & 0.734 & 0.688 & 0.787 \\
推理 & 0.722 & 0.731 & 0.708 & 0.757 \\
ヒューマンドラマ & 0.683 & 0.724 & 0.641 & 0.831 \\
ハイファンタジー & 0.679 & 0.714 & 0.643 & 0.803 \\
現実世界（恋愛） & 0.687 & 0.713 & 0.657 & 0.779 \\
空想科学 & 0.681 & 0.712 & 0.648 & 0.790 \\
純文学 & 0.685 & 0.698 & 0.669 & 0.730 \\
ローファンタジー & 0.646 & 0.694 & 0.611 & 0.803 \\
VRゲーム & 0.596 & 0.682 & 0.561 & 0.868 \\
歴史 & 0.685 & 0.681 & 0.690 & 0.673 \\
ノンジャンル & 0.637 & 0.681 & 0.608 & 0.773 \\
宇宙 & 0.670 & 0.677 & 0.664 & 0.690 \\
詩 & 0.663 & 0.665 & 0.659 & 0.671 \\
その他 & 0.651 & 0.664 & 0.640 & 0.689 \\
パニック & 0.670 & 0.662 & 0.676 & 0.650 \\
エッセイ & 0.593 & 0.620 & 0.581 & 0.665 \\
\end{longtable}


\subsection{サンプリングデータにおける各モデルの学習結果}
LIMEによる解釈性を担保するため，全作品（All）データセットから「エタり」「完結」を各5,000件（計10,000件）抽出し，均衡データセットを作成した．
この同一データセットを用いて，TF-IDF，Bi-LSTM，BERTの3つのモデルで再学習および評価を行った結果を表\ref{tab:text_models_compare}に示す．

頻度ベースのTF-IDF（Accuracy 0.708）や，系列情報を考慮したBi-LSTM（Accuracy 0.744）と比較して，事前学習済みモデルであるBERTはAccuracy 0.899，F1 Score 0.899という圧倒的に高い予測性能を示した．
これは，均衡データを用いることでクラス不均衡の影響が解消されたことに加え，BERTが持つ深い文脈理解能力が，小説のあらすじに含まれる「未完結になりやすい微細なニュアンス」を捉えるのに極めて有効であることを示している．

\begin{table}[H]
  \centering
  \caption{テキストモデルの予測性能比較（均衡データセット: 全作品）}
  \label{tab:text_models_compare}
  \begin{tabular}{lrrrr} \hline
    Model & Accuracy & Precision & Recall & F1 Score \\ \hline
    TF-IDF & 0.71 & 0.71 & 0.71 & 0.70 \\
    Bi-LSTM & 0.74 & 0.75 & 0.74 & 0.76 \\
    \textbf{BERT} & \textbf{0.90} & \textbf{0.90} & \textbf{0.90} & \textbf{0.90} \\ \hline
  \end{tabular}
\end{table}

\section{LIMEによる判断根拠の可視化}
構築したモデルがテキスト中の「どの単語」に着目して判定を行っているかを明らかにするため，LIMEを用いた重要単語の抽出を行った．
以下に，モデル間およびジャンルごとの抽出結果を示す．
なお，各図において，正の値（右側）は「エタり判定」に，負の値（左側）は「完結判定」に寄与した単語を表す．

また，紙面の都合上本節には掲載しきれなかった，全ジャンルにおけるTF-IDFモデルの寄与単語画像については，前述のオンライン補足資料ウェブサイトにて公開している．

\subsection{モデル構造による着目点の差異（全作品対象）}
全作品（All）データを対象に，3つのモデルが重視した単語の傾向を比較する．

\subsubsection{TF-IDFモデル}
TF-IDFモデル（図\ref{fig:lime_all_tfidf}）では，「カクヨム」「転載」「企画」といった，物語の内容そのものよりも作品の掲載状況やメタ情報に関する単語が上位に抽出された．
特筆すべき点として，「完結」という単語がエタり判定（正の方向）に強く寄与していることが挙げられる．これは，「第一部完結」や「本編完結」と記載された状態で続編が更新されず，システム上は未完結（エタり）として扱われているケースを，モデルが特徴として捉えていると解釈できる．

\begin{figure}[H]
  \centering
  \includegraphics[width=0.8\linewidth]{images/lime/all-tfidf.png}
  \caption{LIMEによる判断根拠：全作品・TF-IDFモデル}
  \label{fig:lime_all_tfidf}
\end{figure}

\subsubsection{Bi-LSTMモデル}
Bi-LSTMモデル（図\ref{fig:lime_all_lstm}）では，「……」や「──」といった三点リーダーやダッシュなどの記号，および「だった」などの文末表現が判定根拠として上位に現れた．
TF-IDFと比較して，単語の意味内容だけでなく，記号の使用頻度や文体といったテキストの構造的な特徴を学習し，判定に利用していることがうかがえる．

\begin{figure}[H]
  \centering
    \includegraphics[width=0.8\linewidth]{images/lime/all-lstm.png}
    \caption{LIMEによる判断根拠：全作品・Bi-LSTMモデル}
  \label{fig:lime_all_lstm}
\end{figure}

\subsubsection{BERTモデル}
BERTモデル（図\ref{fig:lime_all_bert}）では，「カク」「\#\#ヨ」（カクヨム）のようにサブワード化されたトークンや，「ます」「です」といった助動詞への着目が見られた．
特に，完結判定（負の方向）に「ます」「です」などの丁寧語（敬体）が寄与している点は特徴的である．これはBERTが，単語レベルの意味だけでなく，あらすじの文体や語り口（丁寧な説明口調か，断定的な口調か等）という高次の文脈情報を捉え，高精度な予測に繋げていることを示唆している．

\begin{figure}[H]
  \centering
  \includegraphics[width=0.8\linewidth]{images/lime/all-bert.png}
  \caption{LIMEによる判断根拠：全作品・BERTモデル}
  \label{fig:lime_all_bert}
\end{figure}

\subsection{ジャンルごとの重要単語（TF-IDF）}
TF-IDFモデルを用いて抽出された，主要ジャンルにおける特徴的な単語を示す．

\subsubsection{異世界（恋愛）}
「異世界（恋愛）」ジャンル（図\ref{fig:lime_isekai_love}）では，「婚約」「ざま（ざまぁ）」といったジャンル特有の展開を示す単語に加え，「カクヨム」「転載」「サイト」など，他サイトでの活動や重複投稿を示唆するメタ情報がエタり判定に強く寄与していた．
これは，作者が他サイトへ活動拠点を移したことによる更新停止や，マルチポスト作品特有の更新傾向を，モデルが学習している可能性を示唆している．

\begin{figure}[H]
  \centering
  \includegraphics[width=0.8\linewidth]{images/lime/101-tfidf.png} 
  \caption{重要単語の可視化：異世界（恋愛）}
  \label{fig:lime_isekai_love}
\end{figure}

\subsubsection{リプレイ}
予測精度が高かった「リプレイ」ジャンル（図\ref{fig:lime_replay}）では，エタり判定に「ダブルクロス」「オーヴァード」「UGN」といった特定のTRPGシステムに関連する固有名称が強く寄与していた．
これは，特定システムのキャンペーン（連続したセッション）リプレイが多数投稿されているものの，長編化しやすく完結に至る前に更新が途絶えるケースが多いことを示唆している．
対照的に，完結判定には「主人公」「物語」「人生」といったストーリー性を強調する一般的な語彙が寄与しており，システムデータやログの羅列ではなく，一つの読み物として構成された作品が完結しやすい傾向が見て取れる．

\begin{figure}[H]
  \centering
  \includegraphics[width=0.8\linewidth]{images/lime/9904-tfidf.png}
  \caption{重要単語の可視化：リプレイ}
  \label{fig:lime_replay}
\end{figure}

\subsubsection{純文学・エッセイ}
「純文学」や「エッセイ」（図\ref{fig:lime_lit_essay}）では，他の物語ジャンルとは異なる興味深い傾向が確認された．
エタり判定には，「連作」「企画」（純文学）や「続編」「参加」（エッセイ）といった，シリーズ化やイベント参加などに関連する単語が寄与している．これは，単発作品よりも継続的な執筆が求められる形式において，維持が困難になるケースが多いことを示唆している．

一方，完結判定には「考え」「綴る」「日記」といった内省的な語彙に加え，両ジャンル共通して「不定期」が強く寄与していた．
ストーリー性の強いジャンルでは「不定期更新＝エタりの前兆」であることが一般的だが，構成の制約が緩いこれらのジャンルにおいては，更新義務に縛られずマイペースに執筆することが，逆に完結（または継続的な活動）に繋がっている可能性を示している．

\begin{figure}[H]
  \centering
  % 上段：純文学
  \begin{subfigure}[b]{1.0\textwidth}
    \centering
    \includegraphics[width=\linewidth]{images/lime/301-tfidf.png}
    \caption{純文学}
  \end{subfigure}
  
  \vspace{1.0cm}
  
  % 下段：エッセイ
  \begin{subfigure}[b]{1.0\textwidth}
    \centering
    \includegraphics[width=\linewidth]{images/lime/9903-tfidf.png}
    \caption{エッセイ}
  \end{subfigure}
  
  \caption{重要単語の可視化：文芸・その他系ジャンル}
  \label{fig:lime_lit_essay}
\end{figure}



% 第5章 考察

\chapter{考察}

本章では，第4章で得られた実験結果に基づき，予測モデルの性能差が生じた要因や，LIMEによって抽出された特徴語が示唆する「連載放棄（エタり）」の構造的なメカニズムについて考察する．

\section{テキスト表現手法による予測精度の差異}
あらすじテキストを用いた実験，特に\textbf{各クラスのデータ数を揃えた均衡データセットにおける比較}において，BERTモデル（Accuracy 0.90）は，TF-IDF（Accuracy 0.71）やBi-LSTM（Accuracy 0.74）と比較して圧倒的に高い性能を示した．この性能差は，各モデルがテキストから抽出できる情報の「質」の違いに起因すると考えられる．

TF-IDFは単語の出現頻度に基づく手法であるため，LIMEの分析結果（図\ref{fig:lime_all_tfidf}）が示すように，「カクヨム」「転載」といった特定のキーワード（メタ情報）に強く依存した予測を行っていた．これは，キーワードが含まれる作品に対しては有効であるが，一般的な語彙で構成された作品の判別は困難であることを意味する．

一方，BERTモデルのLIME結果（図\ref{fig:lime_all_bert}）では，「ます」「です」といった助動詞や，機能語への着目が見られた．これは，BERTが単語の意味だけでなく，文体（敬体か常体か）や，作者の語り口といった高次の文脈情報を捉えていることを示唆している．
例えば，あらすじにおいて読者に丁寧語で語りかけるような作者は，作品の管理や完結に対する意識が高い（あるいは，完結の見通しが立ってから投稿している）といった潜在的な相関関係を，BERTが学習した可能性がある．このように，表面的なキーワードだけでなく，文脈に含まれる微細なニュアンスを考慮できたことが，BERTの高精度な予測に繋がったと推察される．

\section{LIMEによる判断根拠に基づく要因分析}

\subsection{「完結」という単語が示すエタりリスクのパラドックス}
TF-IDFモデルの全作品分析において，「完結」という単語が「エタり判定」に最も強く寄与するという，直感に反する結果が得られた．
あらすじに「完結」と書かれているにも関わらず，システム上のステータスが「未完結（エタり）」となっているこの現象には，主に以下の3つの要因が考えられる．

\begin{enumerate}
    \item \textbf{システム設定の失念}: 
    本文の投稿は完了し，あらすじ等で「完結しました」と報告しているものの，サイトの管理画面での「完結設定」への切り替えを忘れているケース．これは単純なヒューマンエラーであるが，データセット上は「長期間更新がない未完結作品」として扱われる．
    \item \textbf{シリーズ作品の構造的要因}: 
    「第一部完結」「本編完結（番外編へ続く）」といった文脈で使用されているケース．この場合，物語の一区切りはついているものの，続編や番外編の更新が途絶えた時点で，作品全体としてはエタり扱いとなる．
    \item \textbf{商業化に伴う更新停止}: 
    作品が書籍化された際，Web版の更新を停止し，あらすじに「書籍版完結」「書籍化に伴いWeb版はここまで」と記載するケース．これは商業的には成功であるが，Web連載としては放棄された状態となる．
\end{enumerate}

この結果は，テキストマイニングによる予測において，単語の文字通りの意味（完結＝終わっている）と，実際の作品ステータス（連載状況）との間に乖離が生じうることを示唆しており，メタデータとテキスト情報を組み合わせた多角的な判断の必要性を裏付けている．

\subsection{外部プラットフォームと執筆環境の影響}
「異世界（恋愛）」ジャンル等において，「カクヨム」「サイト」「転載」といった単語がエタり判定に強く寄与した．
これは，複数の投稿サイトで作品を公開する「マルチポスト」や，活動拠点の移行が，結果として更新停止のリスクを高めていることを示している．
複数のサイトで並行して連載を維持することは作者にとって大きな負担であり，メインの活動場所を他サイトへ移した場合，元のサイト（小説家になろう）の作品は放置されやすい．LIMEはこの傾向を「エタりの予兆」として的確に捉えていると言える．

\subsection{ジャンル特有の執筆構造と継続性}

\subsubsection*{リプレイにおけるシステム依存リスク}
「リプレイ」ジャンルでは，「ダブルクロス」等の特定TRPGシステム名がエタり要因となり，「物語」等の一般語が完結要因となる傾向が見られた．
特定のシステム名を冠する作品は，実際のゲームセッションをログ化した長編キャンペーンであることが多く，参加者のスケジュール調整やログの編集作業など，完結までに要するコストが非常に高い．
対して，あらすじで「物語」や「人生」を強調する作品は，システムに依存しない創作読み物（リプレイ風小説）や単発セッションである可能性が高く，作者個人のペースで完結させやすいため，このような差異が生まれたと考えられる．

\subsubsection*{非物語ジャンルにおける更新義務の逆説}
「純文学」や「エッセイ」において，「不定期」という単語が完結（継続）に寄与し，「企画」「連作」がエタりに寄与するという，物語ジャンルとは逆の傾向が確認された．
ストーリー漫画や長編小説では，不定期更新は読者離れやエタりの前兆と捉えられがちである．しかし，エッセイや日記的な作品においては，「不定期で書く」という宣言は，作者が更新義務（プレッシャー）から解放され，マイペースに執筆を続けるための防衛策として機能していると考えられる．
逆に，「企画」や「連作」は「期間内に書かなければならない」「オチをつけなければならない」という制約を生むため，モチベーション維持のハードルを上げ，結果として連載放棄に繋がりやすいと推察される．


% 第6章 結論


\chapter{結論}

本研究では，ウェブ小説投稿サイトにおける作品の未完結（エタり）問題を解決するための第一歩として，作品のメタデータおよびあらすじテキストを用いた連載継続予測モデルの構築と評価を行った．

\section{研究の総括}
実験および考察により，以下の知見が得られた．

\begin{enumerate}
    \item \textbf{事前学習済みモデルの有効性}: 
    あらすじを用いた予測において，BERTを用いた手法（正解率0.717）は，TF-IDFやBi-LSTMを用いた手法よりも安定して高い性能を示した．特に，データ数が極端に少ない「リプレイ」ジャンルにおいて，BERTは転移学習の効果により高い再現率を達成し，ロングテールなデータ分布を持つウェブ小説分析において，事前学習済みモデルが極めて有効であることが確認された．
    
    \item \textbf{ジャンル特性と予測可能性の相関}: 
    「異世界（恋愛）」のようにテンプレート性が高く，定型的な表現が多用されるジャンルほど予測性能が高くなる傾向が確認された．一方で，「エッセイ」などの物語構造を持たないジャンルや，「VRゲーム」のように専門用語がノイズとなりやすいジャンルでは，テキスト情報のみからの予測は困難であることが明らかになった．
    
    \item \textbf{メタデータとテキスト情報の比較}: 
    テキスト情報単体でも約72\%の正解率を達成したが，ブックマーク数や評価ポイントなどのメタデータを用いたLightGBMモデル（正解率約78\%）には及ばなかった．これは，作品の内容そのもの以上に，読者からの反応や更新頻度といった外部要因・行動履歴が，執筆継続に強く影響していることを定量的に裏付けている．

    \item \textbf{判断根拠の可視化によるエタり要因の特定}: 
    LIMEを用いた分析により，ジャンルごとにエタりを誘発する特定のキーワードや文脈が明らかになった．特に，「完結」という単語がシステム上のエタり判定と強く相関するパラドックス（完結設定の失念やシリーズ途中での放棄）や，エッセイにおける「不定期」が逆に継続を示唆する傾向など，単なる予測精度を超えた，ウェブ小説特有の執筆実態を浮き彫りにした．
\end{enumerate}

\section{本研究の貢献}
本研究の学術的および社会的貢献は以下の2点に集約される．

第一に，\textbf{大規模なウェブ小説データセットの構築と基礎分析}である．
特定の一時点だけでなく，長期的な更新状況を追跡したデータセットを構築し，ジャンルごとの「エタり率」を定量的に明らかにした．これにより，「童話」や「推理」ジャンルの完結難易度が極めて高いことなど，従来は経験則で語られてきた傾向を数値的根拠に基づいて示した点は，今後のウェブ文学研究における基礎資料として意義がある．

第二に，\textbf{執筆継続性における「内容」と「環境」の影響度の解明}である．
予測モデルの構築およびXAIによる判断根拠の可視化を通じ，執筆を継続するためには，作品の内容（テキスト）もさることながら，更新習慣やマルチポストなどの投稿環境（メタデータ）が極めて重要なファクターであることを実証した．これは，ウェブ小説プラットフォームが作者支援を行う際，「面白い書き方」を教えるだけでなく，「モチベーションを維持する環境作り」や「システム上の完結処理のサポート」が重要であることを示唆している．

\section{今後の展望}
本研究に残された課題と，今後の展望について述べる．

第一に，\textbf{適合率の向上と誤検知の抑制}である．
本研究のBERTモデルは，連載放棄作品を見逃さない「再現率」においては優秀な成績を収めたが，完結する作品を誤ってエタりと判定する「誤検知」も一定数存在した．執筆支援システムとして実用化する場合，過剰な警告は作者の意欲を削ぐ恐れがある．今後は，誤判定に対するコストを考慮した損失関数の設計や，アンサンブル学習によるロバスト性の向上が求められる．

第二に，\textbf{分析結果に基づく具体的な執筆支援システムの開発}である．
本研究ではLIMEによりエタり判定の根拠を可視化することに成功した．今後は，この知見を応用し，執筆中の作者に対してリアルタイムで「エタりリスク」を提示し，具体的な改善案（例：あらすじ内の特定の表現の見直しや，完結設定の確認喚起）をフィードバックする実用的なインターフェースの構築が求められる．

第三に，\textbf{マルチモーダル学習と生成AIによる支援への拡張}である．
現在はメタデータとテキストを別々のモデルで評価しているが，これらを統合したマルチモーダル学習を行うことで，さらなる精度向上が見込まれる．
さらに，単に「エタりそうです」と警告するだけでなく，生成AI（LLM）を用いて「続きの展開案」や「プロットの修正案」を提示することで，連載放棄を未然に防ぐ能動的な執筆支援システムの構築を目指したい．

本研究の成果が，ウェブ小説作者の執筆環境改善に寄与し，一つでも多くの作品が完結を迎え，読者に届けられる一助となることを期待する．



% 謝辞・参考文献


\chapter*{謝辞}
本研究を進めるにあたり，終始懇切丁寧なご指導をいただいた久野雅樹教授に深く感謝いたします．
また，多角的な視点から有益な助言をくださった，同研究室の先輩方に心より感謝申し上げます．

\begin{thebibliography}{99}
  \bibitem{shimizu}
  清水一憲, 伊東栄典, 廣川佐千男, ``集合知に基づくオンライン小説のランキング手法の提案と評価,'' 情報処理学会研究報告, Vol.2013-DBS-157, No.8, pp.1--8, 2013.

  \bibitem{takada}
  高田叶子, 佐藤哲司, ``文体の類似度を考慮したオンライン小説推薦手法の提案,'' 第9回データ工学と情報マネジメントに関するフォーラム (DEIM Forum 2017), B5-2, 2017.

  \bibitem{sakai}
  堺雄之介, 伊東栄典, ``オンライン小説の流行語抽出,'' 情報処理学会第82回全国大会講演論文集, pp.6T-01, 2020.

  \bibitem{honda}
  本田優也, 伊東栄典, ``利用者投稿型小説サイトにおけるキーワードの多様性動向分析,'' 情報処理学会研究報告, Vol.2017-MPS-112, No.4, 2017.

  \bibitem{huang}
  黄晨雯, ``ビッグデータとしてのウェブ小説―言語特徴およびトレンド解析―,'' 大阪大学博士論文, 2022.

  \bibitem{lightgbm}
  G. Ke, Q. Meng, T. Finley, T. Wang, W. Chen, W. Ma, Q. Ye and T. Liu, ``LightGBM: A Highly Efficient Gradient Boosting Decision Tree,'' \textit{Proc. of NIPS}, pp.3146--3154, 2017.

  \bibitem{xgboost}
  T. Chen and C. Guestrin, ``XGBoost: A Scalable Tree Boosting System,'' \textit{Proc. of KDD}, pp.785--794, 2016.

  \bibitem{lstm}
  S. Hochreiter and J. Schmidhuber, ``Long Short-Term Memory,'' \textit{Neural Computation}, Vol.9, No.8, pp.1735--1780, 1997.

  \bibitem{mecab}
  T. Kudo, K. Yamamoto and Y. Matsumoto, ``Applying Conditional Random Fields to Japanese Morphological Analysis,'' \textit{Proc. of EMNLP}, pp.230--237, 2004.

  \bibitem{sentencepiece}
  T. Kudo and J. Richardson, ``SentencePiece: A simple and language independent subword tokenizer and detokenizer for Neural Text Processing,'' \textit{Proc. of EMNLP}, pp.66--71, 2018.

  \bibitem{bert}
  J. Devlin, et al., ``BERT: Pre-training of Deep Bidirectional Transformers for Language Understanding,'' \textit{Proc. of NAACL-HLT}, pp.4171--4186, 2019.

  \bibitem{narou_help}
  ヒナプロジェクト, ``ジャンルについて - 小説家になろうヘルプセンター,'' \url{https://syosetu.com/helpcenter/helppage/helppageid/59}, (参照 2026-01-29).
\end{thebibliography}


% 付録 (Appendix)

\appendix


% 開発環境


\chapter{使用ライブラリ}
\label{app:environment}

本研究では，タスクの性質に応じて複数のPython仮想環境を構築し，実験を行った．
主要なライブラリおよびそのバージョンを表\ref{tab:lib_versions}に示す．

\begin{table}[H]
  \centering
  \small
  \caption{本研究で使用した主要ライブラリとバージョン}
  \label{tab:lib_versions}
  \begin{tabular}{l|l|l} \hline
    ライブラリ名 & バージョン & 主な用途 \\ \hline \hline
    \multicolumn{3}{c}{\textbf{共通・数値分析環境 (torch311)}} \\ \hline
    Python & 3.11.13 & プログラミング言語 \\
    LightGBM & 4.6.0 & 勾配ブースティング決定木 \\
    XGBoost & 3.1.1 & 勾配ブースティング決定木 \\
    scikit-learn & 1.6.1 & 機械学習全般・評価指標 \\
    PyTorch & 2.5.1 & ニューラルネットワーク構築 \\
    MeCab & 0.996.13 & 形態素解析 \\
    SentencePiece & 0.2.1 & サブワード分割 \\
    NumPy & 2.0.1 & 数値計算 \\
    Pandas & 2.3.1 & データフレーム操作 \\ \hline
    \multicolumn{3}{c}{\textbf{大規模言語モデル環境 (torch-bert-env)}} \\ \hline
    Transformers & 4.41.2 & BERTモデルの利用 \\
    Datasets & 4.1.1 & データセット管理 \\
    Accelerate & 1.10.1 & 学習の高速化・最適化 \\
    PyTorch & 2.5.1+cu121 & GPU対応深層学習フレームワーク \\ \hline
    \multicolumn{3}{c}{\textbf{解釈・可視化環境 (lime\_env)}} \\ \hline
    Captum & 0.8.0 & モデル判断根拠の可視化 \\
    LIME & 0.2.0.1 & 局所的な説明可能性 \\
    Transformers & 4.57.1 & 言語モデルの読み込み \\
    PyTorch & 2.7.1+cu118 & 深層学習フレームワーク \\
    Matplotlib & 3.10.7 & グラフ描画 \\
    Seaborn & 0.13.2 & 統計的データ可視化 \\ \hline
  \end{tabular}
\end{table}


% 数値データの詳細


\chapter{数値データに基づく予測性能詳細}
\label{app:numeric_detail}

本付録では，メタデータ（数値特徴量）を用いた機械学習モデルの予測性能について，全ジャンルの詳細な結果を示す．
各表の指標は，正解率（Acc），適合率（Prec），再現率（Rec），F1値（F1）であり，F1値の降順で掲載している．

% --- レイアウト設定 ---
\setlength{\tabcolsep}{3pt} 

% 全データ (単独表示)
\begin{table}[H]
  \centering
  \small
  \caption{全データ (All)}
  \label{tab:res_all}
  \begin{tabular}{l|c|c|c|c} \hline
    Model & Acc & Prec & Rec & F1 \\ \hline
    XGBoost & 0.783 & 0.789 & 0.775 & 0.782 \\
    LightGBM & 0.782 & 0.788 & 0.773 & 0.780 \\
    RandomForest & 0.764 & 0.770 & 0.755 & 0.762 \\
    MLP & 0.754 & 0.749 & 0.765 & 0.756 \\
    SVM & 0.700 & 0.673 & 0.778 & 0.722 \\
    LogisticReg & 0.658 & 0.625 & 0.787 & 0.697 \\
    k-NN & 0.679 & 0.679 & 0.678 & 0.679 \\ \hline
  \end{tabular}
\end{table}

\begin{table}[H]
  \small
  \begin{minipage}[t]{0.49\textwidth}
    \centering
    \caption{異世界（恋愛）}
    \begin{tabular}{l|c|c|c|c} \hline
      Model & Acc & Prec & Rec & F1 \\ \hline
      XGBoost & 0.875 & 0.880 & 0.868 & 0.874 \\
      LightGBM & 0.873 & 0.880 & 0.865 & 0.872 \\
      RandomForest & 0.846 & 0.853 & 0.837 & 0.845 \\
      MLP & 0.841 & 0.859 & 0.817 & 0.837 \\
      SVM & 0.771 & 0.819 & 0.695 & 0.752 \\
      k-NN & 0.745 & 0.748 & 0.739 & 0.744 \\
      LogisticReg & 0.704 & 0.696 & 0.722 & 0.709 \\ \hline
    \end{tabular}
  \end{minipage}
  \hfill
  \begin{minipage}[t]{0.49\textwidth}
    \centering
    \caption{現実世界（恋愛）}
    \begin{tabular}{l|c|c|c|c} \hline
      Model & Acc & Prec & Rec & F1 \\ \hline
      LightGBM & 0.784 & 0.790 & 0.773 & 0.781 \\
      XGBoost & 0.783 & 0.788 & 0.773 & 0.780 \\
      MLP & 0.763 & 0.751 & 0.786 & 0.768 \\
      RandomForest & 0.766 & 0.773 & 0.753 & 0.763 \\
      SVM & 0.713 & 0.712 & 0.718 & 0.715 \\
      k-NN & 0.684 & 0.684 & 0.684 & 0.684 \\
      LogisticReg & 0.628 & 0.624 & 0.646 & 0.634 \\ \hline
    \end{tabular}
  \end{minipage}
\end{table}

\begin{table}[H]
  \small
  \begin{minipage}[t]{0.49\textwidth}
    \centering
    \caption{ハイファンタジー}
    \begin{tabular}{l|c|c|c|c} \hline
      Model & Acc & Prec & Rec & F1 \\ \hline
      LightGBM & 0.789 & 0.789 & 0.788 & 0.788 \\
      XGBoost & 0.788 & 0.789 & 0.787 & 0.788 \\
      RandomForest & 0.768 & 0.768 & 0.768 & 0.768 \\
      MLP & 0.753 & 0.755 & 0.749 & 0.752 \\
      SVM & 0.714 & 0.704 & 0.741 & 0.722 \\
      k-NN & 0.683 & 0.683 & 0.684 & 0.684 \\
      LogisticReg & 0.619 & 0.610 & 0.660 & 0.634 \\ \hline
    \end{tabular}
  \end{minipage}
  \hfill
  \begin{minipage}[t]{0.49\textwidth}
    \centering
    \caption{ローファンタジー}
    \begin{tabular}{l|c|c|c|c} \hline
      Model & Acc & Prec & Rec & F1 \\ \hline
      XGBoost & 0.764 & 0.762 & 0.768 & 0.765 \\
      LightGBM & 0.765 & 0.765 & 0.764 & 0.765 \\
      MLP & 0.737 & 0.730 & 0.754 & 0.742 \\
      RandomForest & 0.741 & 0.744 & 0.736 & 0.740 \\
      SVM & 0.692 & 0.669 & 0.757 & 0.710 \\
      k-NN & 0.677 & 0.674 & 0.686 & 0.680 \\
      LogisticReg & 0.631 & 0.624 & 0.658 & 0.641 \\ \hline
    \end{tabular}
  \end{minipage}
\end{table}

\begin{table}[H]
  \small
  \begin{minipage}[t]{0.49\textwidth}
    \centering
    \caption{純文学}
    \begin{tabular}{l|c|c|c|c} \hline
      Model & Acc & Prec & Rec & F1 \\ \hline
      XGBoost & 0.783 & 0.773 & 0.801 & 0.787 \\
      LightGBM & 0.778 & 0.776 & 0.783 & 0.779 \\
      RandomForest & 0.772 & 0.777 & 0.763 & 0.770 \\
      MLP & 0.757 & 0.745 & 0.781 & 0.763 \\
      SVM & 0.733 & 0.695 & 0.832 & 0.757 \\
      k-NN & 0.709 & 0.695 & 0.749 & 0.721 \\
      LogisticReg & 0.688 & 0.673 & 0.734 & 0.702 \\ \hline
    \end{tabular}
  \end{minipage}
  \hfill
  \begin{minipage}[t]{0.49\textwidth}
    \centering
    \caption{ヒューマンドラマ}
    \begin{tabular}{l|c|c|c|c} \hline
      Model & Acc & Prec & Rec & F1 \\ \hline
      XGBoost & 0.766 & 0.773 & 0.755 & 0.763 \\
      LightGBM & 0.767 & 0.777 & 0.750 & 0.763 \\
      MLP & 0.754 & 0.756 & 0.750 & 0.753 \\
      RandomForest & 0.755 & 0.767 & 0.734 & 0.750 \\
      SVM & 0.705 & 0.665 & 0.829 & 0.738 \\
      k-NN & 0.694 & 0.689 & 0.708 & 0.698 \\
      LogisticReg & 0.666 & 0.639 & 0.765 & 0.696 \\ \hline
    \end{tabular}
  \end{minipage}
\end{table}

\begin{table}[H]
  \small
  \begin{minipage}[t]{0.49\textwidth}
    \centering
    \caption{歴史}
    \begin{tabular}{l|c|c|c|c} \hline
      Model & Acc & Prec & Rec & F1 \\ \hline
      LightGBM & 0.770 & 0.768 & 0.775 & 0.771 \\
      XGBoost & 0.767 & 0.763 & 0.775 & 0.769 \\
      MLP & 0.756 & 0.742 & 0.787 & 0.763 \\
      RandomForest & 0.754 & 0.760 & 0.742 & 0.751 \\
      SVM & 0.714 & 0.679 & 0.812 & 0.740 \\
      LogisticReg & 0.684 & 0.664 & 0.743 & 0.701 \\
      k-NN & 0.689 & 0.679 & 0.716 & 0.697 \\ \hline
    \end{tabular}
  \end{minipage}
  \hfill
  \begin{minipage}[t]{0.49\textwidth}
    \centering
    \caption{推理}
    \begin{tabular}{l|c|c|c|c} \hline
      Model & Acc & Prec & Rec & F1 \\ \hline
      XGBoost & 0.790 & 0.796 & 0.780 & 0.788 \\
      LightGBM & 0.783 & 0.792 & 0.766 & 0.779 \\
      RandomForest & 0.780 & 0.790 & 0.764 & 0.777 \\
      MLP & 0.762 & 0.759 & 0.769 & 0.764 \\
      SVM & 0.712 & 0.687 & 0.778 & 0.729 \\
      k-NN & 0.705 & 0.696 & 0.727 & 0.711 \\
      LogisticReg & 0.667 & 0.653 & 0.710 & 0.680 \\ \hline
    \end{tabular}
  \end{minipage}
\end{table}

\begin{table}[H]
  \small
  \begin{minipage}[t]{0.49\textwidth}
    \centering
    \caption{ホラー}
    \begin{tabular}{l|c|c|c|c} \hline
      Model & Acc & Prec & Rec & F1 \\ \hline
      LightGBM & 0.783 & 0.775 & 0.798 & 0.786 \\
      XGBoost & 0.783 & 0.782 & 0.785 & 0.783 \\
      MLP & 0.773 & 0.763 & 0.795 & 0.778 \\
      RandomForest & 0.779 & 0.790 & 0.762 & 0.775 \\
      SVM & 0.735 & 0.684 & 0.874 & 0.767 \\
      LogisticReg & 0.718 & 0.687 & 0.803 & 0.740 \\
      k-NN & 0.684 & 0.664 & 0.748 & 0.703 \\ \hline
    \end{tabular}
  \end{minipage}
  \hfill
  \begin{minipage}[t]{0.49\textwidth}
    \centering
    \caption{アクション}
    \begin{tabular}{l|c|c|c|c} \hline
      Model & Acc & Prec & Rec & F1 \\ \hline
      XGBoost & 0.761 & 0.772 & 0.741 & 0.756 \\
      LightGBM & 0.758 & 0.767 & 0.742 & 0.754 \\
      RandomForest & 0.754 & 0.769 & 0.726 & 0.747 \\
      MLP & 0.746 & 0.756 & 0.729 & 0.742 \\
      SVM & 0.719 & 0.740 & 0.676 & 0.706 \\
      k-NN & 0.688 & 0.692 & 0.677 & 0.684 \\
      LogisticReg & 0.653 & 0.635 & 0.719 & 0.674 \\ \hline
    \end{tabular}
  \end{minipage}
\end{table}

\begin{table}[H]
  \small
  \begin{minipage}[t]{0.49\textwidth}
    \centering
    \caption{コメディー}
    \begin{tabular}{l|c|c|c|c} \hline
      Model & Acc & Prec & Rec & F1 \\ \hline
      XGBoost & 0.737 & 0.745 & 0.722 & 0.733 \\
      LightGBM & 0.733 & 0.739 & 0.720 & 0.730 \\
      RandomForest & 0.730 & 0.736 & 0.719 & 0.727 \\
      MLP & 0.724 & 0.729 & 0.714 & 0.721 \\
      SVM & 0.687 & 0.687 & 0.689 & 0.687 \\
      k-NN & 0.658 & 0.659 & 0.657 & 0.658 \\
      LogisticReg & 0.632 & 0.615 & 0.701 & 0.656 \\ \hline
    \end{tabular}
  \end{minipage}
  \hfill
  \begin{minipage}[t]{0.49\textwidth}
    \centering
    \caption{VRゲーム}
    \begin{tabular}{l|c|c|c|c} \hline
      Model & Acc & Prec & Rec & F1 \\ \hline
      XGBoost & 0.734 & 0.739 & 0.723 & 0.731 \\
      LightGBM & 0.730 & 0.736 & 0.718 & 0.726 \\
      RandomForest & 0.721 & 0.726 & 0.711 & 0.718 \\
      MLP & 0.705 & 0.706 & 0.703 & 0.704 \\
      SVM & 0.654 & 0.647 & 0.678 & 0.662 \\
      LogisticReg & 0.624 & 0.617 & 0.654 & 0.635 \\
      k-NN & 0.616 & 0.615 & 0.623 & 0.619 \\ \hline
    \end{tabular}
  \end{minipage}
\end{table}

\begin{table}[H]
  \small
  \begin{minipage}[t]{0.49\textwidth}
    \centering
    \caption{宇宙}
    \begin{tabular}{l|c|c|c|c} \hline
      Model & Acc & Prec & Rec & F1 \\ \hline
      LightGBM & 0.739 & 0.734 & 0.749 & 0.741 \\
      XGBoost & 0.740 & 0.736 & 0.746 & 0.741 \\
      MLP & 0.714 & 0.703 & 0.742 & 0.721 \\
      RandomForest & 0.720 & 0.721 & 0.720 & 0.720 \\
      SVM & 0.693 & 0.673 & 0.750 & 0.709 \\
      k-NN & 0.687 & 0.683 & 0.698 & 0.690 \\
      LogisticReg & 0.671 & 0.664 & 0.690 & 0.677 \\ \hline
    \end{tabular}
  \end{minipage}
  \hfill
  \begin{minipage}[t]{0.49\textwidth}
    \centering
    \caption{空想科学}
    \begin{tabular}{l|c|c|c|c} \hline
      Model & Acc & Prec & Rec & F1 \\ \hline
      XGBoost & 0.774 & 0.767 & 0.787 & 0.777 \\
      LightGBM & 0.773 & 0.767 & 0.785 & 0.775 \\
      MLP & 0.760 & 0.745 & 0.792 & 0.768 \\
      RandomForest & 0.764 & 0.763 & 0.767 & 0.765 \\
      SVM & 0.714 & 0.687 & 0.788 & 0.733 \\
      k-NN & 0.701 & 0.694 & 0.718 & 0.706 \\
      LogisticReg & 0.667 & 0.661 & 0.688 & 0.674 \\ \hline
    \end{tabular}
  \end{minipage}
\end{table}

\begin{table}[H]
  \small
  \begin{minipage}[t]{0.49\textwidth}
    \centering
    \caption{パニック}
    \begin{tabular}{l|c|c|c|c} \hline
      Model & Acc & Prec & Rec & F1 \\ \hline
      LightGBM & 0.732 & 0.739 & 0.719 & 0.728 \\
      XGBoost & 0.725 & 0.731 & 0.714 & 0.722 \\
      RandomForest & 0.728 & 0.742 & 0.700 & 0.720 \\
      SVM & 0.670 & 0.639 & 0.781 & 0.703 \\
      MLP & 0.691 & 0.678 & 0.727 & 0.700 \\
      LogisticReg & 0.665 & 0.643 & 0.743 & 0.689 \\
      k-NN & 0.659 & 0.656 & 0.671 & 0.663 \\ \hline
    \end{tabular}
  \end{minipage}
  \hfill
  \begin{minipage}[t]{0.49\textwidth}
    \centering
    \caption{童話}
    \begin{tabular}{l|c|c|c|c} \hline
      Model & Acc & Prec & Rec & F1 \\ \hline
      MLP & 0.789 & 0.792 & 0.791 & 0.789 \\
      LightGBM & 0.779 & 0.785 & 0.774 & 0.777 \\
      RandomForest & 0.784 & 0.809 & 0.748 & 0.775 \\
      XGBoost & 0.774 & 0.780 & 0.764 & 0.772 \\
      SVM & 0.749 & 0.720 & 0.824 & 0.768 \\
      LogisticReg & 0.746 & 0.720 & 0.807 & 0.760 \\
      k-NN & 0.739 & 0.724 & 0.781 & 0.750 \\ \hline
    \end{tabular}
  \end{minipage}
\end{table}

\begin{table}[H]
  \small
  \begin{minipage}[t]{0.49\textwidth}
    \centering
    \caption{詩}
    \begin{tabular}{l|c|c|c|c} \hline
      Model & Acc & Prec & Rec & F1 \\ \hline
      XGBoost & 0.737 & 0.734 & 0.745 & 0.738 \\
      RandomForest & 0.742 & 0.750 & 0.728 & 0.738 \\
      LightGBM & 0.735 & 0.736 & 0.735 & 0.735 \\
      MLP & 0.719 & 0.714 & 0.737 & 0.723 \\
      SVM & 0.692 & 0.668 & 0.764 & 0.711 \\
      LogisticReg & 0.671 & 0.652 & 0.728 & 0.687 \\
      k-NN & 0.679 & 0.671 & 0.701 & 0.685 \\ \hline
    \end{tabular}
  \end{minipage}
  \hfill
  \begin{minipage}[t]{0.49\textwidth}
    \centering
    \caption{エッセイ}
    \begin{tabular}{l|c|c|c|c} \hline
      Model & Acc & Prec & Rec & F1 \\ \hline
      XGBoost & 0.746 & 0.741 & 0.756 & 0.748 \\
      LightGBM & 0.743 & 0.738 & 0.753 & 0.745 \\
      MLP & 0.734 & 0.723 & 0.761 & 0.741 \\
      SVM & 0.715 & 0.684 & 0.802 & 0.738 \\
      RandomForest & 0.736 & 0.738 & 0.731 & 0.735 \\
      LogisticReg & 0.695 & 0.675 & 0.750 & 0.711 \\
      k-NN & 0.684 & 0.681 & 0.694 & 0.687 \\ \hline
    \end{tabular}
  \end{minipage}
\end{table}

\begin{table}[H]
  \small
  \begin{minipage}[t]{0.49\textwidth}
    \centering
    \caption{リプレイ}
    \begin{tabular}{l|c|c|c|c} \hline
      Model & Acc & Prec & Rec & F1 \\ \hline
      RandomForest & 0.766 & 0.762 & 0.779 & 0.768 \\
      XGBoost & 0.754 & 0.742 & 0.792 & 0.762 \\
      LightGBM & 0.729 & 0.708 & 0.779 & 0.739 \\
      MLP & 0.716 & 0.715 & 0.754 & 0.728 \\
      SVM & 0.698 & 0.687 & 0.754 & 0.714 \\
      k-NN & 0.667 & 0.634 & 0.803 & 0.705 \\
      LogisticReg & 0.654 & 0.655 & 0.692 & 0.665 \\ \hline
    \end{tabular}
  \end{minipage}
  \hfill
  \begin{minipage}[t]{0.49\textwidth}
    \centering
    \caption{その他}
    \begin{tabular}{l|c|c|c|c} \hline
      Model & Acc & Prec & Rec & F1 \\ \hline
      XGBoost & 0.751 & 0.748 & 0.758 & 0.753 \\
      LightGBM & 0.750 & 0.751 & 0.747 & 0.749 \\
      MLP & 0.739 & 0.723 & 0.778 & 0.749 \\
      RandomForest & 0.741 & 0.745 & 0.734 & 0.740 \\
      SVM & 0.692 & 0.660 & 0.792 & 0.720 \\
      k-NN & 0.689 & 0.685 & 0.700 & 0.692 \\
      LogisticReg & 0.654 & 0.639 & 0.708 & 0.672 \\ \hline
    \end{tabular}
  \end{minipage}
\end{table}

\begin{table}[H]
  \small
  \begin{minipage}[t]{0.49\textwidth}
    \centering
    \caption{ノンジャンル}
    \begin{tabular}{l|c|c|c|c} \hline
      Model & Acc & Prec & Rec & F1 \\ \hline
      LightGBM & 0.775 & 0.793 & 0.745 & 0.768 \\
      XGBoost & 0.775 & 0.794 & 0.742 & 0.767 \\
      RandomForest & 0.758 & 0.775 & 0.727 & 0.750 \\
      MLP & 0.742 & 0.759 & 0.709 & 0.733 \\
      SVM & 0.681 & 0.679 & 0.688 & 0.684 \\
      LogisticReg & 0.663 & 0.657 & 0.683 & 0.670 \\
      k-NN & 0.648 & 0.650 & 0.642 & 0.646 \\ \hline
    \end{tabular}
  \end{minipage}
  \hfill
  \begin{minipage}[t]{0.49\textwidth}
  \end{minipage}
\end{table}



% TF-IDF詳細


\chapter{TF-IDFモデルのハイパーパラメータ探索結果}
\label{app:tfidf_grid}

TF-IDFとロジスティック回帰を用いた分類における，グリッドサーチ結果の詳細を以下に示す．
各表は，特徴量次元数（Feat）と正則化パラメータ（$C$）を変化させた際の検証性能（正解率：Accuracy）である．
太字は各ジャンルにおける最高性能を示す．

% 全データ
\begin{table}[H]
    \centering
    \small
    \caption{全データ (All)}
    \begin{tabular}{l|rrrrr}
        \hline
        Feat \textbackslash C & 0.05 & 0.08 & 0.1 & 0.2 & 0.5 \\ \hline
        18000 & 0.6836 & 0.6874 & 0.6896 & 0.6911 & 0.6908 \\
        20000 & 0.6839 & 0.6877 & 0.6894 & \textbf{0.6919} & 0.6916 \\
        22000 & 0.6841 & 0.6875 & 0.6893 & 0.6918 & 0.6909 \\ \hline
    \end{tabular}
\end{table}

\begin{table}[H]
    \centering
    \small
    \begin{minipage}[t]{0.48\textwidth}
        \centering
        \caption{異世界（恋愛）}
        \begin{tabular}{l|rrrrr}
            \hline
            Feat \textbackslash C & 0.5 & 0.8 & 1.0 & 1.2 & 1.5 \\ \hline
            8000 & 0.7378 & \textbf{0.7393} & 0.7371 & 0.7369 & 0.7350 \\
            10000 & 0.7373 & 0.7387 & 0.7376 & 0.7363 & 0.7344 \\
            12000 & 0.7363 & 0.7376 & 0.7377 & 0.7373 & 0.7348 \\ \hline
        \end{tabular}
    \end{minipage}
    \hfill
    \begin{minipage}[t]{0.48\textwidth}
        \centering
        \caption{現実世界（恋愛）}
        \begin{tabular}{l|rrrrr}
            \hline
            Feat \textbackslash C & 0.5 & 0.8 & 1.0 & 1.2 & 1.5 \\ \hline
            18000 & 0.6791 & 0.6808 & 0.6806 & 0.6801 & 0.6777 \\
            20000 & 0.6786 & 0.6814 & 0.6821 & 0.6801 & 0.6786 \\
            22000 & 0.6791 & \textbf{0.6843} & 0.6838 & 0.6828 & 0.6805 \\ \hline
        \end{tabular}
    \end{minipage}
\end{table}

\begin{table}[H]
    \centering
    \small
    \begin{minipage}[t]{0.48\textwidth}
        \centering
        \caption{ハイファンタジー}
        \begin{tabular}{l|rrrrr}
            \hline
            Feat \textbackslash C & 0.5 & 0.8 & 1.0 & 1.2 & 1.5 \\ \hline
            18000 & 0.6684 & 0.6660 & 0.6656 & 0.6646 & 0.6640 \\
            20000 & \textbf{0.6697} & 0.6671 & 0.6662 & 0.6668 & 0.6654 \\
            22000 & 0.6696 & 0.6686 & 0.6678 & 0.6675 & 0.6664 \\ \hline
        \end{tabular}
    \end{minipage}
    \hfill
    \begin{minipage}[t]{0.48\textwidth}
        \centering
        \caption{ローファンタジー}
        \begin{tabular}{l|rrrrr}
            \hline
            Feat \textbackslash C & 0.5 & 0.8 & 1.0 & 1.2 & 1.5 \\ \hline
            3000 & \textbf{0.6417} & 0.6405 & 0.6399 & 0.6389 & 0.6381 \\
            5000 & 0.6383 & 0.6356 & 0.6340 & 0.6338 & 0.6291 \\
            7000 & 0.6389 & 0.6372 & 0.6372 & 0.6364 & 0.6342 \\ \hline
        \end{tabular}
    \end{minipage}
\end{table}

\begin{table}[H]
    \centering
    \small
    \begin{minipage}[t]{0.48\textwidth}
        \centering
        \caption{純文学}
        \begin{tabular}{l|rrrrr}
            \hline
            Feat \textbackslash C & 0.5 & 0.8 & 1.0 & 1.2 & 1.5 \\ \hline
            800 & 0.6215 & 0.6257 & 0.6240 & 0.6232 & 0.6198 \\
            1000 & 0.6215 & \textbf{0.6266} & 0.6249 & 0.6249 & 0.6206 \\
            1200 & 0.6189 & 0.6215 & 0.6198 & 0.6240 & 0.6206 \\ \hline
        \end{tabular}
    \end{minipage}
    \hfill
    \begin{minipage}[t]{0.48\textwidth}
        \centering
        \caption{ヒューマンドラマ}
        \begin{tabular}{l|rrrrr}
            \hline
            Feat \textbackslash C & 0.5 & 0.8 & 1.0 & 1.2 & 1.5 \\ \hline
            18000 & \textbf{0.6849} & 0.6793 & 0.6807 & 0.6803 & 0.6774 \\
            20000 & 0.6836 & 0.6807 & 0.6805 & 0.6818 & 0.6783 \\
            22000 & 0.6820 & 0.6811 & 0.6797 & 0.6813 & 0.6774 \\ \hline
        \end{tabular}
    \end{minipage}
\end{table}

\begin{table}[H]
    \centering
    \small
    \begin{minipage}[t]{0.48\textwidth}
        \centering
        \caption{歴史}
        \begin{tabular}{l|rrrrr}
            \hline
            Feat \textbackslash C & 0.5 & 0.8 & 1.0 & 1.2 & 1.5 \\ \hline
            2000 & 0.6476 & 0.6438 & 0.6418 & 0.6389 & 0.6399 \\
            5000 & 0.6505 & 0.6534 & 0.6525 & 0.6525 & 0.6525 \\
            7000 & \textbf{0.6554} & 0.6544 & 0.6486 & 0.6438 & 0.6457 \\ \hline
        \end{tabular}
    \end{minipage}
    \hfill
    \begin{minipage}[t]{0.48\textwidth}
        \centering
        \caption{推理}
        \begin{tabular}{l|rrrrr}
            \hline
            Feat \textbackslash C & 0.5 & 0.8 & 1.0 & 1.2 & 1.5 \\ \hline
            18000 & 0.6806 & 0.6827 & \textbf{0.6848} & 0.6806 & 0.6784 \\
            20000 & 0.6806 & 0.6827 & \textbf{0.6848} & 0.6806 & 0.6784 \\
            22000 & 0.6806 & 0.6827 & \textbf{0.6848} & 0.6806 & 0.6784 \\ \hline
        \end{tabular}
    \end{minipage}
\end{table}

\begin{table}[H]
    \centering
    \small
    \begin{minipage}[t]{0.48\textwidth}
        \centering
        \caption{ホラー}
        \begin{tabular}{l|rrrrr}
            \hline
            Feat \textbackslash C & 0.5 & 0.8 & 1.0 & 1.2 & 1.5 \\ \hline
            8000 & 0.6719 & 0.6661 & 0.6701 & 0.6649 & 0.6603 \\
            10000 & 0.6719 & 0.6730 & \textbf{0.6736} & 0.6713 & 0.6661 \\
            12000 & 0.6707 & 0.6701 & 0.6724 & 0.6701 & 0.6655 \\ \hline
        \end{tabular}
    \end{minipage}
    \hfill
    \begin{minipage}[t]{0.48\textwidth}
        \centering
        \caption{アクション}
        \begin{tabular}{l|rrrrr}
            \hline
            Feat \textbackslash C & 0.5 & 0.8 & 1.0 & 1.2 & 1.5 \\ \hline
            800 & 0.6682 & 0.6674 & 0.6713 & 0.6667 & 0.6667 \\
            1000 & 0.6644 & 0.6713 & \textbf{0.6797} & 0.6782 & 0.6720 \\
            1200 & 0.6674 & 0.6705 & 0.6790 & 0.6767 & 0.6743 \\ \hline
        \end{tabular}
    \end{minipage}
\end{table}

\begin{table}[H]
    \centering
    \small
    \begin{minipage}[t]{0.48\textwidth}
        \centering
        \caption{コメディー}
        \begin{tabular}{l|rrrrr}
            \hline
            Feat \textbackslash C & 0.5 & 0.8 & 1.0 & 1.2 & 1.5 \\ \hline
            18000 & 0.6678 & 0.6678 & \textbf{0.6700} & 0.6667 & 0.6667 \\
            20000 & 0.6678 & 0.6678 & \textbf{0.6700} & 0.6667 & 0.6667 \\
            22000 & 0.6678 & 0.6678 & \textbf{0.6700} & 0.6667 & 0.6667 \\ \hline
        \end{tabular}
    \end{minipage}
    \hfill
    \begin{minipage}[t]{0.48\textwidth}
        \centering
        \caption{VRゲーム}
        \begin{tabular}{l|rrrrr}
            \hline
            Feat \textbackslash C & 0.05 & 0.08 & 0.1 & 0.12 & 0.5 \\ \hline
            800 & 0.6263 & \textbf{0.6386} & \textbf{0.6386} & \textbf{0.6386} & 0.6283 \\
            1000 & 0.6304 & 0.6283 & 0.6345 & 0.6366 & 0.6263 \\
            1200 & 0.6263 & 0.6222 & 0.6263 & 0.6242 & 0.6263 \\ \hline
        \end{tabular}
    \end{minipage}
\end{table}

\begin{table}[H]
    \centering
    \small
    \begin{minipage}[t]{0.48\textwidth}
        \centering
        \caption{宇宙}
        \begin{tabular}{l|rrrrr}
            \hline
            Feat \textbackslash C & 0.5 & 0.8 & 1.0 & 1.2 & 2.0 \\ \hline
            800 & 0.6528 & 0.6572 & 0.6507 & 0.6485 & 0.6376 \\
            1000 & 0.6594 & 0.6638 & 0.6638 & \textbf{0.6681} & 0.6594 \\
            1200 & 0.6419 & 0.6485 & 0.6441 & 0.6485 & 0.6528 \\ \hline
        \end{tabular}
    \end{minipage}
    \hfill
    \begin{minipage}[t]{0.48\textwidth}
        \centering
        \caption{空想科学}
        \begin{tabular}{l|rrrrr}
            \hline
            Feat \textbackslash C & 0.5 & 0.8 & 1.0 & 1.2 & 2.0 \\ \hline
            3000 & 0.6350 & 0.6423 & 0.6423 & 0.6437 & 0.6470 \\
            5000 & 0.6430 & 0.6496 & \textbf{0.6516} & 0.6470 & 0.6496 \\
            7000 & 0.6284 & 0.6357 & 0.6397 & 0.6417 & 0.6377 \\ \hline
        \end{tabular}
    \end{minipage}
\end{table}

\begin{table}[H]
    \centering
    \small
    \begin{minipage}[t]{0.48\textwidth}
        \centering
        \caption{パニック}
        \begin{tabular}{l|rrrrr}
            \hline
            Feat \textbackslash C & 0.05 & 0.08 & 0.1 & 0.2 & 0.5 \\ \hline
            700 & 0.6635 & 0.6667 & \textbf{0.6730} & 0.6571 & 0.6698 \\
            1000 & \textbf{0.6730} & 0.6667 & 0.6698 & 0.6667 & 0.6508 \\
            2000 & 0.6698 & 0.6667 & 0.6667 & 0.6571 & 0.6508 \\ \hline
        \end{tabular}
    \end{minipage}
    \hfill
    \begin{minipage}[t]{0.48\textwidth}
        \centering
        \caption{童話}
        \begin{tabular}{l|rrrrr}
            \hline
            Feat \textbackslash C & 0.05 & 0.08 & 0.1 & 0.2 & 0.5 \\ \hline
            3000 & 0.6737 & \textbf{0.6780} & \textbf{0.6780} & \textbf{0.6780} & 0.6674 \\
            5000 & 0.6737 & \textbf{0.6780} & \textbf{0.6780} & \textbf{0.6780} & 0.6674 \\
            7000 & 0.6737 & \textbf{0.6780} & \textbf{0.6780} & \textbf{0.6780} & 0.6674 \\ \hline
        \end{tabular}
    \end{minipage}
\end{table}

\begin{table}[H]
    \centering
    \small
    \begin{minipage}[t]{0.48\textwidth}
        \centering
        \caption{詩}
        \begin{tabular}{l|rrrrr}
            \hline
            Feat \textbackslash C & 0.05 & 0.08 & 0.1 & 0.2 & 0.5 \\ \hline
            700 & 0.6369 & 0.6311 & 0.6167 & 0.6023 & 0.6110 \\
            1000 & \textbf{0.6513} & 0.6455 & 0.6427 & 0.6282 & 0.6311 \\
            2000 & 0.6455 & 0.6398 & 0.6427 & 0.6398 & 0.6311 \\ \hline
        \end{tabular}
    \end{minipage}
    \hfill
    \begin{minipage}[t]{0.48\textwidth}
        \centering
        \caption{エッセイ}
        \begin{tabular}{l|rrrrr}
            \hline
            Feat \textbackslash C & 0.5 & 0.8 & 1.0 & 1.5 & 2.0 \\ \hline
            700 & 0.6345 & 0.6354 & 0.6297 & 0.6288 & 0.6269 \\
            1000 & 0.6458 & 0.6458 & 0.6430 & 0.6326 & 0.6269 \\
            2000 & 0.6430 & 0.6420 & 0.6449 & \textbf{0.6468} & 0.6402 \\ \hline
        \end{tabular}
    \end{minipage}
\end{table}

\begin{table}[H]
    \centering
    \small
    \begin{minipage}[t]{0.48\textwidth}
        \centering
        \caption{リプレイ}
        \begin{tabular}{l|rrrrr}
            \hline
            Feat \textbackslash C & 0.05 & 0.08 & 0.1 & 0.2 & 0.5 \\ \hline
            700 & \textbf{0.8154} & 0.8000 & 0.8000 & \textbf{0.8154} & 0.8000 \\
            1000 & \textbf{0.8154} & 0.8000 & 0.8000 & \textbf{0.8154} & 0.8000 \\
            2000 & \textbf{0.8154} & 0.8000 & 0.8000 & \textbf{0.8154} & 0.8000 \\ \hline
        \end{tabular}
    \end{minipage}
    \hfill
    \begin{minipage}[t]{0.48\textwidth}
        \centering
        \caption{その他}
        \begin{tabular}{l|rrrrr}
            \hline
            Feat \textbackslash C & 0.5 & 0.8 & 1.0 & 1.2 & 2.0 \\ \hline
            8000 & 0.6274 & 0.6262 & 0.6274 & 0.6250 & 0.6208 \\
            10000 & 0.6256 & 0.6250 & \textbf{0.6286} & \textbf{0.6286} & 0.6220 \\
            12000 & 0.6256 & 0.6250 & \textbf{0.6286} & \textbf{0.6286} & 0.6220 \\ \hline
        \end{tabular}
    \end{minipage}
\end{table}

\begin{table}[H]
    \centering
    \small
    \begin{minipage}[t]{0.48\textwidth}
        \centering
        \caption{ノンジャンル}
        \begin{tabular}{l|rrrrr}
            \hline
            Feat \textbackslash C & 0.05 & 0.08 & 0.1 & 0.2 & 0.5 \\ \hline
            3000 & 0.6291 & 0.6356 & 0.6366 & 0.6375 & 0.6356 \\
            5000 & 0.6318 & 0.6368 & 0.6399 & \textbf{0.6430} & 0.6379 \\
            7000 & 0.6287 & 0.6344 & 0.6384 & 0.6429 & 0.6388 \\ \hline
        \end{tabular}
    \end{minipage}
    \hfill
    \begin{minipage}[t]{0.48\textwidth}
    \end{minipage}
\end{table}


% Bi-LSTM詳細

\chapter{Bi-LSTMモデルのクラス別評価指標}
\label{app:bilstm_detail}

Bi-LSTMモデルにおける，各ジャンルの適合率（Precision），再現率（Recall），F1値（F1-score）の詳細な内訳（Classification Report）を以下に示す．

\begin{longtable}{lrrr}
\caption{Bi-LSTM Classification Report Summary} \\
\hline
Genre / Class & Precision & Recall & F1-score \\ \hline
\endfirsthead

\multicolumn{4}{c}{{続・Bi-LSTM Classification Report Summary}} \\
\hline
Genre / Class & Precision & Recall & F1-score \\ \hline
\endhead

\hline \multicolumn{4}{r}{\textit{Continued on next page}} \\
\endfoot

\hline
\endlastfoot

% --- データ本体 ---
\multicolumn{4}{l}{\textbf{全データ (All)}} \\
\quad 非エタり & 0.70 & 0.73 & 0.71 \\
\quad エタり & 0.72 & 0.68 & 0.70 \\ \hline

\multicolumn{4}{l}{\textbf{異世界（恋愛）}} \\
\quad 非エタり & 0.76 & 0.66 & 0.70 \\
\quad エタり & 0.70 & 0.79 & 0.74 \\ \hline

\multicolumn{4}{l}{\textbf{現実世界（恋愛）}} \\
\quad 非エタり & 0.64 & 0.77 & 0.70 \\
\quad エタり & 0.71 & 0.57 & 0.63 \\ \hline

\multicolumn{4}{l}{\textbf{ハイファンタジー}} \\
\quad 非エタり & 0.65 & 0.68 & 0.66 \\
\quad エタり & 0.66 & 0.63 & 0.65 \\ \hline

\multicolumn{4}{l}{\textbf{ローファンタジー}} \\
\quad 非エタり & 0.62 & 0.66 & 0.64 \\
\quad エタり & 0.64 & 0.60 & 0.62 \\ \hline

\multicolumn{4}{l}{\textbf{純文学}} \\
\quad 非エタり & 0.59 & 0.53 & 0.56 \\
\quad エタり & 0.57 & 0.63 & 0.60 \\ \hline

\multicolumn{4}{l}{\textbf{ヒューマンドラマ}} \\
\quad 非エタり & 0.67 & 0.71 & 0.69 \\
\quad エタり & 0.69 & 0.65 & 0.67 \\ \hline

\multicolumn{4}{l}{\textbf{歴史}} \\
\quad 非エタり & 0.63 & 0.60 & 0.61 \\
\quad エタり & 0.61 & 0.64 & 0.63 \\ \hline

\multicolumn{4}{l}{\textbf{推理}} \\
\quad 非エタり & 0.64 & 0.71 & 0.68 \\
\quad エタり & 0.68 & 0.60 & 0.64 \\ \hline

\multicolumn{4}{l}{\textbf{ホラー}} \\
\quad 非エタり & 0.62 & 0.73 & 0.67 \\
\quad エタり & 0.67 & 0.55 & 0.60 \\ \hline

\multicolumn{4}{l}{\textbf{アクション}} \\
\quad 非エタり & 0.64 & 0.74 & 0.68 \\
\quad エタり & 0.69 & 0.58 & 0.63 \\ \hline

\multicolumn{4}{l}{\textbf{コメディー}} \\
\quad 非エタり & 0.62 & 0.67 & 0.64 \\
\quad エタり & 0.64 & 0.58 & 0.61 \\ \hline

\multicolumn{4}{l}{\textbf{VRゲーム}} \\
\quad 非エタり & 0.58 & 0.65 & 0.61 \\
\quad エタり & 0.60 & 0.53 & 0.57 \\ \hline

\multicolumn{4}{l}{\textbf{宇宙}} \\
\quad 非エタり & 0.62 & 0.55 & 0.58 \\
\quad エタり & 0.59 & 0.66 & 0.62 \\ \hline

\multicolumn{4}{l}{\textbf{空想科学}} \\
\quad 非エタり & 0.58 & 0.73 & 0.65 \\
\quad エタり & 0.64 & 0.47 & 0.54 \\ \hline

\multicolumn{4}{l}{\textbf{パニック}} \\
\quad 非エタり & 0.62 & 0.55 & 0.58 \\
\quad エタり & 0.59 & 0.66 & 0.63 \\ \hline

\multicolumn{4}{l}{\textbf{童話}} \\
\quad 非エタり & 0.67 & 0.66 & 0.66 \\
\quad エタり & 0.66 & 0.68 & 0.67 \\ \hline

\multicolumn{4}{l}{\textbf{詩}} \\
\quad 非エタり & 0.57 & 0.84 & 0.68 \\
\quad エタり & 0.69 & 0.36 & 0.47 \\ \hline

\multicolumn{4}{l}{\textbf{エッセイ}} \\
\quad 非エタり & 0.59 & 0.66 & 0.62 \\
\quad エタり & 0.61 & 0.53 & 0.57 \\ \hline

\multicolumn{4}{l}{\textbf{リプレイ} (予測失敗事例)} \\
\quad 非エタり & 0.51 & 1.00 & 0.67 \\
\quad エタり & \textbf{0.00} & \textbf{0.00} & \textbf{0.00} \\ \hline

\multicolumn{4}{l}{\textbf{その他}} \\
\quad 非エタり & 0.58 & 0.73 & 0.64 \\
\quad エタり & 0.63 & 0.46 & 0.53 \\ \hline

\multicolumn{4}{l}{\textbf{ノンジャンル}} \\
\quad 非エタり & 0.61 & 0.74 & 0.67 \\
\quad エタり & 0.67 & 0.53 & 0.59 \\ \hline

\end{longtable}



% オンライン補足資料

\chapter{オンライン補足資料}
\label{app:online_supplement}

本論文内で掲載しきれなかった詳細な実験結果画像については，以下の特設ウェブサイトにて公開している．
サイト内では，全ジャンルにおける「LightGBM，XGBoost，RandomForestそれぞれの特徴量重要度」「誤分類データの密度分布」「Bi-LSTMモデルの学習曲線（Loss推移）」を閲覧可能である．


\begin{itemize}
    \item \textbf{サイト名}: 令和7年度 卒業論文 付録資料 
    \item \textbf{URL}: \url{https://sites.google.com/gl.cc.uec.ac.jp/sato-2025-research?usp=sharing}
\end{itemize}

\end{document}